\documentclass[11pt,oneside,openany]{book}
\usepackage[T1]{fontenc}
\usepackage[utf8]{inputenc}
\usepackage{lmodern,amsmath,amssymb,amsthm,mathtools,booktabs,longtable}
\usepackage[margin=1in,headheight=15pt]{geometry}
\usepackage{microtype,graphicx,tikz,tikz-cd,fancyhdr}
\pagestyle{fancy}
\fancyhf{}
\fancyhead[L]{\small Surface Orders and Cyclic Time}
\fancyhead[R]{\small\thechapter}
\fancyfoot[C]{\thepage}
\usetikzlibrary{arrows.meta,positioning,calc,decorations.markings}
\usepackage[hidelinks]{hyperref}
\hypersetup{pdftitle={Surface Orders, Cyclic Time and the Arithmetic Spectral Bridge},pdfauthor={Amelie Schreiber}}
\newtheorem{theorem}{Theorem}[chapter]
\newtheorem{proposition}[theorem]{Proposition}
\newtheorem{lemma}[theorem]{Lemma}
\newtheorem{corollary}[theorem]{Corollary}
\theoremstyle{definition}\newtheorem{definition}[theorem]{Definition}
\newtheorem{example}[theorem]{Example}
\theoremstyle{remark}\newtheorem{remark}[theorem]{Remark}
\newcommand{\C}{\mathbb C}\newcommand{\Z}{\mathbb Z}\newcommand{\R}{\mathbb R}
\newcommand{\gshift}[2]{#1\langle #2\rangle}
\DeclareMathOperator{\tr}{tr}\DeclareMathOperator{\Rea}{Re}
\newcommand{\status}{\textbf{Research book. GRH is not proved here.}}
\setcounter{tocdepth}{1}
\title{Surface Orders, Cyclic Time\ and the Arithmetic Spectral Bridge\\[1em]\large Dessins, resolutions, Auslander--Reiten theory, quantum models and verifiable research}
\author{Amelie Schreiber\\\small Research development edition}
\date{November 5, 2026}
\begin{document}
\chapter{Validated Dirichlet targets and finite readout certificates}
% Generated shared reviewed material.
\section{A finite, validated approximation result}
Let $\chi_3=(0,1,-1)$, $\chi_4=(0,1,0,-1)$ and
$\chi_5=(0,1,-1,-1,1)$, extended periodically, with parities $a=1,1,0$.
We approximate the completed functions
\[
 \Lambda_q(s)=(q/\pi)^{(s+a)/2}\Gamma((s+a)/2)L(s,\chi_q).
\]
The primitive theta/Mellin identity is the explicit classical analytic input;
no hypothesis about zeros is used. These three examples do not assert a
uniform result for all conductors, general Hecke functions or Artin holomorphy.

\begin{theorem}[Scope of the supplied finite certificate]
For each $q=3,4,5$, the supplied exact-rational calculation constructs a
complex polynomial $P_q$ of degree $28$ satisfying
\[
 \sup_{|s-1/2|\le7/10}|P_q(s)-\Lambda_q(s)|<5\cdot10^{-15}.
\]
The computation includes enclosures of all $64$ true Fourier-circle nodes,
completed-function sample error, analytic truncation, Fourier aliasing,
the exact closed-form $180$-step gradient-descent iterate and coefficient
midpoint rounding. The Python implementation and informal analytic proof
have independent checks; this is not a kernel-checked Lean theorem.
\end{theorem}
\begin{proof}
Write $F(c+z)=\sum_{n\ge0}a_nz^n$, with $c=1/2$, $R=7/10$,
$\rho=6/5$, $r=12/5$. The independently recomputed rational theta majorants
are strictly below $M_q=4,10,6$, respectively, on $|z|\le r$.
Use true roots $\zeta_j=e^{2\pi i j/64}$ and rational labels $y_j$
with proved error $|y_j-F(c+\rho\zeta_j)|\le\eta$.
Use the real Euclidean gradient on complex coefficient coordinates for
\[
 \mathcal L(b)=\frac1{128}\sum_{j=0}^{63}
 \left|\sum_{n=0}^{28}b_n\zeta_j^n-y_j\right|^2.
\]
Discrete Fourier orthogonality gives the least-squares optimum
$b_n^*=64^{-1}\sum_jy_j\bar\zeta_j^n$, $0\le n\le28$,
and gradient $\nabla\mathcal L(b)=b-b^*$.
For step $1/4$ and zero initialization, the exact $180$-step iterate is
$b_n^{180}=(1-(3/4)^{180})b_n^*$. The code evaluates this closed form;
it does not claim an explicit180-update optimizer run.

Set $u=R/\rho=7/12$, $v=R/r=7/24$ and $w=(\rho/r)^{64}=2^{-64}$.
Let $\widehat b_n$ be an outward rectangle enclosing $b_n^*$ and
$\epsilon_n$ the sum of the coordinate halfwidths of the rectangle
for $b_n^{180}$. Choose each actual polynomial coefficient as that
rectangle's exact rational midpoint $c_n$. The normalized basis is explicit:
\[
 P_q(s)=\sum_{n=0}^{28}c_n\left(\frac{s-1/2}{6/5}\right)^n.
\]
Thus the exported coefficient with index $n$ multiplies the normalized
variable, rather than the raw power $(s-1/2)^n$. Its error on the compact disk is bounded by
\begin{align*}
 E_{\rm opt}&=(3/4)^{180}\sum_{n=0}^{28}|\widehat b_n|_1u^n,&
 E_{\rm arithmetic}&=\sum_{n=0}^{28}\epsilon_nu^n,\\
 E_{\rm oracle}&=\eta\frac{1-u^{29}}{1-u},&
 E_{\rm alias}&=M_q\frac{w}{1-w}\frac{1-v^{29}}{1-v},\\
 E_{\rm Taylor}&=M_q\frac{v^{29}}{1-v}.&&
\end{align*}
Here $|\widehat b_n|_1$ is the sum of the maximum absolute real and imaginary
coordinates, a safe upper bound for every enclosed complex modulus.
Cauchy's estimate $|a_n|\le M_qr^{-n}$ and the Fourier identity
$b_n^*=a_n\rho^n+\sum_{\ell\ge1}a_{n+64\ell}\rho^{n+64\ell}$
for exact labels prove the alias and Taylor bounds. Label perturbations
contribute at most $\eta$ per Fourier coefficient. The remaining two
terms follow directly from the exact iterate and coefficient rectangles.
Each term is rounded upward to a $180$-bit dyadic rational; the reported total
is the exact sum of those reported terms. The certified totals below are
strictly below the asserted threshold. Detailed sample enclosures are proved
in the extended edition and encoded with exact integer endpoints.
\end{proof}

% Generated shared reviewed material.
\begin{center}\begin{tabular}{c|cc}Conductor & Maximum sample radius & Uniform polynomial bound\\\hline
3 & $3.026\times10^{-17}$ & $1.786\times10^{-15}$ \\
4 & $7.206\times10^{-17}$ & $4.455\times10^{-15}$ \\
5 & $1.404\times10^{-16}$ & $2.906\times10^{-15}$ \\
\end{tabular}\end{center}
The displayed values are rounded upward; exact rational records establish the strict inequalities.


\section{What this does and does not certify}
The final JSON records exact rational centers, interval endpoints, source
hashes and all five error terms. Float values are display summaries only.
A high-precision comparison is an independent diagnostic, not the reason
an interval contains its mathematical target. Arithmetic uses outward
integer operations, and invalid float/order inputs are rejected, including
through typed caches. Earlier pilots and failed validation probes are retained.

This is a fixed finite tolerance/configuration. Increasing accuracy needs
new truncation and precision choices; the fixed $N=11,T=128$ interface is
not an arbitrary-accuracy oracle. The broader approximation existence theorem
is separate. The certificate says nothing about the location of zeros:
$P_q$ is not required to be real-rooted or a self-adjoint spectral determinant.
The independent arithmetic determinant identity, positivity and original
signed short-row/covariance estimates remain open. No GRH proof or complete
new Lean4 formalization is claimed.

% Generated shared reviewed material.
\section{Validated primitive operations}
Real intervals are closed dyadic intervals $[l2^{-224},h2^{-224}]$ with
integer endpoints; complex intervals are Cartesian products. Addition is
exact at this scale. For multiplication take the minimum/maximum of all four
integer endpoint products, divide by $2^{224}$ and round down/up. For a
reciprocal whose interval omits zero, reverse the endpoint order and round
$2^{448}/h$ down and $2^{448}/l$ up. Negative denominators and scalar division
use the same endpoint rules. Each operation encloses all represented real
values by monotonicity; induction proves inclusion for every expression
built from them. Rectangular complex multiplication uses these real rules.
The sum of coordinate halfwidths bounds the complex distance from its exact
rational center. Binary floating arguments are never silently converted.

For $0<x<1$, the alternating arctangent sum and the endpoint obtained by
adding the next signed term enclose $\arctan x$. Sixty terms at $1/5$ and
$1/239$, together with Machin's identity
$\pi=16\arctan(1/5)-4\arctan(1/239)$, give the interval for $\pi$.
The negative coefficient reverses the second interval. The identity follows
from tangent addition and the range of the two angles; no decimal value of
$\pi$ is a trusted input.

For rational $x>0$, range-reduce $x=2^kv$, $1\le v<2$. With
$z=(v-1)/(v+1)\le1/3$,
\[
 \log v=2\sum_{j=0}^{J-1}\frac{z^{2j+1}}{2j+1}+\mathcal R,
 \qquad 0\le\mathcal R\le
 \frac{2z^{2J+1}}{(2J+1)(1-z^2)}.
\]
The geometric tail bounds each remaining denominator below by $2J+1$.
Use $J=92$, the analogous expression for $\log2$, and sign-aware interval
multiplication by $k$. This covers negative $k$ as well.

For complex $z$, repeatedly halve until its interval $\ell^1$ magnitude
bound is $A\le1/2$. The degree $60$ Taylor polynomial for $e^z$ then has remainder
at most
\[
 \frac{A^{61}}{61!}\frac1{1-A/62}.
\]
Successive absolute terms have ratio at most $A/62$, so geometric comparison
proves this bound. Add its symmetric interval to each coordinate, then square
outward the required number of times. The identity $e^{2z}=(e^z)^2$ and
interval inclusion prove the final enclosure. Caches are type-sensitive:
otherwise a cached exact argument can accidentally bypass rejection of an
equal-valued float. No implementation change is allowed during the recorded run;
the producer verifies source bytes again before emitting its certificate.

\section{Recomputing the compact growth constants}
For the larger growth disk, $c=1/2$, $r=12/5$, take $k=2$ when $a=1$
and $k=1$ when $a=0$. Both real Mellin exponents on that disk are at most
$k$. Define
\[
 Q_k(x)=\sum_{j=0}^k\frac{k!}{(k-j)!x^{j+1}},\qquad
 U_{16}(x)=(1+x/16)^{-16},\qquad \delta=3/q.
\]
Integration by parts gives
$\int_1^\infty t^ke^{-xt}\,dt=e^{-x}Q_k(x)$.
The same Gaussian split and geometric-tail comparison used below, now
with cutoff $n_0=8$, imply the exact rational majorant
\[
 \widehat M_q=2\left[
 \sum_{n=1}^{7}n^aU_{16}(\delta n^2)Q_k(\delta n^2)
 +C_aQ_k(64\delta)\frac{U_{16}(32\delta)}{1-U_{16}(8\delta)}\right],
 \quad C_0=1,\quad C_1=(1+\delta^{-1})/2.
\]
Replacing the actual $d=\pi/q$ by the smaller positive $\delta$ increases
the integral majorant. Exact fraction comparison gives
$\widehat M_3<4$, $\widehat M_4<10$, $\widehat M_5<6$.
The producer recomputes these fractions, asserts the strict comparisons and
exports them; the constants in the readout proof are thus checked rather
than supplied by a floating growth estimate.

\section{The completed-function integral and its tails}
Let $d=\pi/q$ and $\theta_q(t)=\sum_{n\ge1}\chi_q(n)n^ae^{-dn^2t}$.
The standard primitive theta transformation, followed by splitting the Mellin
integral at1, gives
\[
 \Lambda_q(s)=\int_1^\infty\theta_q(t)t^{(s+a)/2-1}\,dt+
              \int_1^\infty\theta_q(t)t^{(1-s+a)/2-1}\,dt.
\]
The second root-number coefficient is1 for these three characters: their
finite Gauss sums are $i\sqrt3$, $2i$ and $\sqrt5$, respectively.
For the last, $2\cos(2\pi/5)-2\cos(4\pi/5)=\sqrt5$; the first two are direct
conjugate-root differences. Division by $i^a\sqrt q$ gives the claimed
normalization. Exponential decay makes the two integrals entire. This classical
theta identity, not an assumption of Artin holomorphy or desired positivity,
is the analytic input of the implementation.

The actual circle nodes are enclosed by computing $e^{2\pi i j/64}$ with
the validated primitives. All input boxes fit
$D=[-1,2]+i[-4/3,4/3]$. Thus both exponents
$p=(s+a)/2-1$ and $p'=(1-s+a)/2-1$ satisfy
$|\Re p|\le2$, $|\Im p|\le2/3$. The domain guard precedes all term discards.
Retain $1\le n\le11$ and integrate $1\le t\le128$.

Put $\delta=3/q<d$, $U(x)=(1+x/128)^{-128}$ and
$Q(x)=x^{-1}+2x^{-2}+2x^{-3}$. Since $(1+x/128)^{128}\le e^x$,
$e^{-x}\le U(x)<1$ for $x>0$. Integration by parts gives
$\int_1^\infty t^2e^{-xt}\,dt=e^{-x}Q(x)$.
With $C_0=1$, $C_1=(1+\delta^{-1})/2$ the omitted-$n$ tail of both integrals is
\[
 2C_aQ(\delta12^2)
 \frac{U(\delta12^2/2)}{1-U(12\delta)}.
\]
Indeed $n^ae^{-\delta n^2/2}\le C_a$, the function $Q$ decreases, and
$(12+j)^2\ge12^2+24j$. These inequalities yield a geometric bound after
splitting the exponent in halves. The retained-$n$ time tail is at most
\[
 \sum_{n=1}^{11}2n^a128^3U(128\delta n^2)Q(128\delta n^2).
\]
It follows by substituting $t=128u$ and majorizing $t^{\Re p}$ by $t^2$.
The two tails are disjoint: omitted $n$ use the full integral, whereas
retained $n$ contribute only $t>128$.

\section{High-order integration with a certified remainder}
Split each dyadic interval $[2^j,2^{j+1}]$, $j=0,\ldots,6$, into four equal
panels. For its midpoint $m$ and halfwidth $h$, $h/m\le1/9$.
For $f(z)=e^{-\lambda z}z^p$, $\lambda=\pi n^2/q$, the principal logarithm
is analytic on $|z-m|\le\rho=m/2$. On that circle,
\[
 |f(z)|\le M_\rho=3\max\{(3m/2)^2,(m/2)^{-2},1\}
 U(\delta n^2m/2).
\]
Here $\Re z\ge m/2$, $|\arg z|\le1$ and
$e^{|\Im p|}\le e^{2/3}\le3$. The maximum handles both signs of $\Re p$.
Cauchy's coefficient estimate and cancellation of odd powers on $[-h,h]$
therefore bound the error after degree $2K$ by
\[
 \frac{2hM_\rho(h/\rho)^{2K+2}}
 {(2K+3)(1-(h/\rho)^2)}.
\]
Indeed integrate each remaining even term, replace its denominator by the
first omitted denominator, then sum a geometric series. Here $K=12$.
Multiply by $2n^a$ for both Mellin terms and sum all $28$ panels and retained $n$.
Keeping this majorant even for later discarded terms only overcounts error.

To evaluate the polynomial stably, set $u=h/m$, $q_{-1}=0$, $q_0=1$ and use
\[
 q_{j+1}=\frac{u(p-\lambda m-j)q_j-(\lambda m)u^2q_{j-1}}{j+1}.
\]
Multiplying $f'=(-\lambda+p/t)f$ by $t=m+x$ and matching coefficients proves
this recurrence, with $q_j=f^{(j)}(m)h^j/(j!f(m))$. The integrated polynomial is
$2he^{-\lambda m}m^p\sum_{j=0}^{K}q_{2j}/(2j+1)$; all its arithmetic is outward.

Fixed absolute exponential-rounding errors can be amplified by high-order
coefficients in late panels. The code therefore first bounds a complete
$(n,\text{panel})$ pair, including both powers, by
\[
 4hn^a(m+h)^2U(\delta n^2(m-h)).
\]
For real $t\ge1$ its phase has modulus1 and $t^{\Re p}\le(m+h)^2$, proving
the bound. If it is below $2^{-90}$, skip that pair entirely and add its bound
to the final radius. This is controlled truncation, not numerical underflow.
The final sample radius is the interval-sum halfwidth plus the omitted-$n$,
time, Cauchy and discard bounds. All are recorded separately and checked as
nonnegative. No high-precision special-function routine is trusted here.

\section{Reproduction and remaining interfaces}
Run the preserved source on the existing Linux Python environment with no new
packages: \texttt{python3} on \path{examples/analytic_oracle/theta_oracle.py}, with argument \texttt{--output NEW}.
It refuses to overwrite an existing experiment directory. The final-v5
experiment contains $192$ samples, three rational degree $28$ polynomials, copied
source and SHA-256 hashes. Independent validation checks rational endpoint
arithmetic, transcendental enclosures, $1160$ exact coefficient identities, high-
precision completed-function comparisons and the exported component sums.

The fixed-parameter oracle solves a concrete evaluation/certification gap.
It does not identify a positive arithmetic operator, force real-rooted
polynomials, prove the signed original-family short-row bound or supply
complete Lean4 verification. Surface/order and neural architectures may use
these certificates as independent targets, but a small learned loss is not a
substitute for the remaining arithmetic and categorical theorems.

% Generated shared reviewed material.
\section{Requested-error readouts from preserved certified labels}\label{ar:requested}
The accepted final-v5 sample certificate is preserved unchanged. For each
$q=3,4,5$, it supplies $m=64$ rational complex labels $y_j$, bounds
$|y_j-\Lambda_q(c+\rho\zeta_j)|\le\eta_j$, and outward rectangles enclosing
the true roots $\zeta_j=e^{2\pi i j/m}$, where $c=1/2$, $\rho=6/5$.
The classical primitive theta/Mellin identity and accepted growth bounds
$M_q=4,10,6$ on $|s-c|\le r=12/5$ remain the analytic inputs.
This extension accepts an exact rational error request $\varepsilon>0$ on
the fixed compact $|s-c|\le R=7/10$. It changes the finite readout, while
retaining the certified sample errors and the $64$-node alias budget.
It is not an arbitrary-accuracy theta-sampling interface.

\begin{theorem}[Executed rational readout with explicit refusal]\label{ar:returned}
For each supported conductor and rational request, the finite-data
selection below either returns a normalized rational polynomial $P$ and an
executed, exactly replayable gradient-descent trace with
\[
 \sup_{|s-c|\le R}|P(s)-\Lambda_q(s)|\le\varepsilon,
\]
or refuses the request with the exact minimum base budget over all
$64$ allowed degrees. Refusal concerns attainability by this particular
fixed-data majorant; it is not an impossibility theorem for approximation.
\end{theorem}
\begin{proof}
Let $B_n$ be an outward rectangle enclosing
$b_n^*=m^{-1}\sum_j y_j\bar\zeta_j^n$, $0\le n<64$, computed from the
preserved root rectangles. Put $t_n=\operatorname{mid}(B_n)$ and let
$e_n$ be the sum of its coordinate halfwidths. These are exact rationals,
and $|t_n-b_n^*|\le e_n$.
For a selected $d<64$ execute, using exact rational operations,
\[
 b_n^0=0,\qquad b_n^{k+1}=\tfrac34 b_n^k+\tfrac14 t_n,
 \qquad 0\le n\le d.
\]
This is genuine gradient descent with step $1/4$ for the explicitly
specified real Euclidean objective
$J(b)=\frac12\sum_{n=0}^d|b_n-t_n|^2$, whose gradient is $b-t$.
All iterates and exact objective values are retained. The output is literally
\[
 P_k(s)=\sum_{n=0}^d b_n^k\left(\frac{s-c}{\rho}\right)^n.
\]
Set $u=R/\rho=7/12$, $v=R/r=7/24$, $w=(\rho/r)^{64}=2^{-64}$,
$S_d(x)=\sum_{n=0}^d x^n$ and $\eta=\max_j\eta_j$. The five budgets are
\begin{align*}
 E_{\rm opt}(k)&=\sum_{n=0}^d|b_n^k-t_n|_1u^n,&
 E_{\rm oracle}&=\eta S_d(u),\\
 E_{\rm alias}&=M_q\frac{w}{1-w}S_d(v),&
 E_{\rm Taylor}&=M_q\frac{v^{d+1}}{1-v},\\
 E_{\rm arithmetic}&=\sum_{n=0}^d e_nu^n.&&
\end{align*}
Here $|x+iy|_1=|x|+|y|$. Cauchy's estimate and true-root Fourier
orthogonality give the alias and Taylor terms; certified label errors give
the oracle term. Triangle inequality applied to
$b_n^k-b_n^*=(b_n^k-t_n)+(t_n-b_n^*)$ gives the other two.
Thus their exact sum bounds the compact-uniform error.

Define $C_d=E_{\rm oracle}+E_{\rm alias}+E_{\rm Taylor}
+E_{\rm arithmetic}$. Select the smallest $d<64$ with $C_d<\varepsilon$;
if none exists, report $\min_{d<64}C_d$ and refuse. For a selected degree,
$E_{\rm opt}(k)=(3/4)^k\sum_n|t_n|_1u^n\longrightarrow0$.
Therefore executing updates until
$E_{\rm opt}(k)\le\varepsilon-C_d$ terminates after finitely many steps.
The serialized coefficients, budgets, targets and trace use exact rational
endpoints; displayed decimal approximations are not proof inputs.
\end{proof}

% Generated shared reviewed material.
\section{Objective semantics, precision floors and replay}\label{ar:semantics}
The original true-circle sample-label least-squares objective is
\[
 \mathcal L_y(b)=\frac1{2m}\sum_{j=0}^{m-1}
 \left|\sum_{n=0}^d b_n\zeta_j^n-y_j\right|^2.
\]
With $d<m$, Fourier orthogonality gives $A^*A=mI$ and the real Euclidean
coordinate gradient $\nabla\mathcal L_y(b)=b-b^*$. The executed objective
$J$ instead has the exact rational midpoint target $t$.
The difference $\nabla\mathcal L_y-\nabla J=t-b^*$ is enclosed
coefficientwise by $e_n$.
Consequently this is an actual rational GD run for $J$ and a certified
perturbed-gradient implementation for $\mathcal L_y$. It is not the exact
optimizer trajectory of the ideal-root objective, nor a replay of neural or
quantum training. An interval enclosure of an ideal-root GD trajectory
would be a different explicitly labelled object.

Each actual update contracts $b^k-t$ by $3/4$, so the exact stored losses
satisfy $J(b^{k+1})=(9/16)J(b^k)$. A replay recomputes every rational
coordinate update, loss and weighted residual. The closed form
$b^k=(1-(3/4)^k)t$ is an independent check on the executed trace, not its
producer. Exact rational loop coordinates introduce no further midpoint or
floating export error. The arithmetic budget covers the difference between
the certified Fourier rectangles and their midpoint targets, including the
inherited root enclosure and outward interval operations. A subsequent
float conversion requires its own exact conversion-error bound.

The normalized polynomial convention is essential: a stored coefficient
with index $n$ multiplies $((s-1/2)/(6/5))^n$, not $(s-1/2)^n$.
The input certificate is pinned by SHA-256, including its source hashes,
all $192$ sample enclosures, root rectangles and analytic radius components.
Producing source is captured at entry, checked unchanged before emission,
and copied into a fresh output directory. An existing output is never
overwritten. An independent replay should also check the fixed input hash,
all dimensions, degree choice, the first stopping iterate, all five exact
terms and their literal sum against the requested rational tolerance.

The attainable boundary of this majorant is
\[
 C_{\min,q}=\min_{0\le d<64}C_{q,d}.
\]
Every request $\varepsilon>C_{\min,q}$ has strict budget slack for some
degree and the proved finite stop rule. At or below that boundary this
interface refuses. The degree minimizing $C_d$ need not be $63$:
Taylor error decreases with degree while the oracle and arithmetic sums
increase. In particular $C_d\ge\eta$ and
$C_d\ge M_q2^{-64}/(1-2^{-64})$. Increasing the number of optimization
steps does not lower these fixed-data budgets. Approximate boundary values
on the stated compact are $7.2929\cdot10^{-17}$,
$1.7370\cdot10^{-16}$ and $3.3741\cdot10^{-16}$ for $q=3,4,5$;
only the exact exported rational minima are used for acceptance.
These are certificate floors, not lower bounds on the actual errors.

A full adaptive theta oracle would separately need proof-preserving choices
of omitted-$n$ cutoff, time cutoff, quadrature order, whole-panel discard
budget, primitive series lengths, interval precision and serialized-bound
precision. The present $N=11$, $T=128$, $224$-bit arithmetic,
$60/92/60$ primitive series lengths, $2^{-90}$ discard threshold and
$180$-bit rounded budgets cannot support every positive request.
Increasing a nominal bit count alone leaves the fixed primitive remainders
and analytic tails in place. The preserved labels are not recomputed here.
The classical identity, finite arithmetic and informal proof are not new
Lean declarations. No zero-location, all-conductor, positive arithmetic
determinant, quantum-advantage or GRH conclusion follows.

The interactive replay is available at \href{https://amelie-iska.github.io/GRH-research/readout.html}{the certified readout laboratory}. All decimal graphics are illustrative views of the linked exact records.

% Generated shared reviewed material.
\section{Requested-error theta samples for three fixed conductors}\label{at:sampling}
For the primitive characters $\chi_3=(0,1,-1)$, $\chi_4=(0,1,0,-1)$,
$\chi_5=(0,1,-1,-1,1)$ with parities $a=1,1,0$, retain the classical
primitive theta/Mellin identity and root number one already stated above.
Put $D=[-1,2]+i[-4/3,4/3]$ and $\delta=3/q$.
The following adaptive extension is implemented with independent interval
and sample checks; it is an informal analytic theorem, not a Lean result.
\begin{theorem}[Finite adaptive sample specification]\label{at:sample}
For every supported conductor, exact rational $s\in D$, and rational
$\eta>0$, the choices below terminate and give a rational complex midpoint
$y$ and exact rational radius $R\le\eta$ with
$|y-\Lambda_q(s)|\le R$.
The same conclusion holds for a specified true circle node when its
root and node enclosures are refined on every precision retry.
\end{theorem}
\begin{proof}
For $x\ge0$ define the exact rational bound $E(x)=2^{-\lfloor x\rfloor}$.
Since $e>2$, $e^{-x}\le E(x)$; this bound is nonincreasing and tends to zero.
Let $Q(x)=x^{-1}+2x^{-2}+2x^{-3}$ and
$C_0=1$, $C_1=(1+\delta^{-1})/2$.
For $N\ge1$, $n_0=N+1$, $J\ge1$, $T=2^J$, use
\begin{align*}
 B_n&=2C_aQ(\delta n_0^2)
       \frac{E(\delta n_0^2/2)}{1-E(\delta n_0)},\\
 B_t&=\sum_{n=1}^N2n^aT^3E(\delta n^2T)Q(\delta n^2T).
\end{align*}
Here $\delta n_0\ge6/5$ ensures a positive denominator; $E(x)$ equals
one for $0\le x<1$, so this safeguard is essential.
On the four panels of each $[2^j,2^{j+1}]$, write midpoint $m$,
halfwidth $h$, $\rho=m/2$, $v=h/\rho\le2/9$,
and $B(m)=\max\{(m+\rho)^2,(m-\rho)^{-2},1\}$.
Use the Cauchy remainder
\[
 B_Q=\sum_{\mathrm{panels}}\sum_{n=1}^N
 \frac{2n^a\,2h\,3B(m)E(\delta n^2m/2)v^{2K+2}}
 {(2K+3)(1-v^2)}.
\]
Choose in turn the first $N\ge1$, first $J\ge1$, and first $K\ge0$
for which their respective budgets are at most $\eta/8$.
Gaussian decay, then exponential time decay, then $v\le2/9$ prove
termination. These are the accepted tail and Cauchy arguments with each
true exponential separately replaced by its upper bound $E$.

There are $P=4JN$ potential pairs. A nonzero-character pair may be
omitted only if its entire two-power contribution has bound
\[
 b=4hn^a(m+h)^2E(\delta n^2(m-h))\le\eta/(8P).
\]
Add every omitted bound to $B_{\rm discard}$, giving
$B_{\rm discard}\le\eta/8$.
Evaluate all other integrated Taylor polynomials by the stated exact
coefficient recurrence using outward rectangles. Refine both interval
precision and all primitive series remainders until the final rectangle's
sum of coordinate halfwidths satisfies $R_{\rm arith}\le\eta/4$.
For this fixed finite expression all input and primitive intervals shrink;
continuity and outward arithmetic imply $R_{\rm arith}\to0$.
Precision doubling therefore terminates, and the sum of the five raw
components is at most $3\eta/4$.
Choose an integer output scale $B$ such that $5\,2^{-B}\le\eta/8$ and
round each component upward by
$\operatorname{up}_B(x)=2^{-B}\lceil2^Bx\rceil$.
The reported radius is their literal exact sum. With
$\mathbf b=(R_{\rm arith},B_n,B_t,B_Q,B_{\rm discard})$, it obeys
\[
 |y-\Lambda_q(s)|\le\sum_{i=1}^5\operatorname{up}_B(b_i)
 \le\frac{3\eta}{4}+5\,2^{-B}\le\frac{7\eta}{8}<\eta.
\]
The output retains exact rational midpoints, rounded components, $B$,
raw interval radius, parameters, input metadata and source hashes.
The raw analytic bounds are exactly reconstructible; they need not be
serialized as potentially enormous fractions. The output precision varies
with the request and introduces no fixed reporting floor.
\end{proof}

This fixed-conductor sample result uses no hypothesis on zeros.
It does not remove a downstream fixed-node Fourier alias floor, supply
an all-conductor theorem or a positive arithmetic determinant, or prove GRH.

% Generated shared reviewed material.
\section{Adaptive sampling parameters and precision termination}\label{at:details}
The primitive theta/Mellin identity with root number one yields
\[
 \Lambda_q(s)=\sum_{n\ge1}\chi_q(n)n^a
 \int_1^\infty e^{-\pi n^2t/q}
 \bigl(t^{(s+a)/2-1}+t^{(1-s+a)/2-1}\bigr)\,dt.
\]
On $D=[-1,2]+i[-4/3,4/3]$, both exponents $p$ satisfy
$|\Re p|\le2$ and $|\Im p|\le2/3$.
For exact $x\ge0$, put $E(x)=2^{-\lfloor x\rfloor}$.
The elementary inequality $e>2$ gives
$e^{-x}\le e^{-\lfloor x\rfloor}\le E(x)$ and
$E(x)\le2e^{-x\log2}$.
Use this solely as a pointwise nonincreasing upper bound; it is not
multiplicative and has no asserted exponential integral identity.

Integration by parts gives
$\int_1^\infty t^2e^{-xt}\,dt=e^{-x}Q(x)$ for
$Q(x)=x^{-1}+2x^{-2}+2x^{-3}$.
For $a=0$ the bound $n^ae^{-\delta n^2/2}\le C_0=1$ is immediate.
For $a=1$, $n\le(1+n^2)/2$ and
$n^2e^{-\delta n^2/2}\le2/(e\delta)\le1/\delta$ give
$C_1=(1+1/\delta)/2$.
For $n=n_0+j$, the inequality
$n^2\ge n_0^2+2n_0j$ therefore bounds the omitted-$n$ sum by
\[
 2C_aQ(\delta n_0^2)
 \frac{e^{-\delta n_0^2/2}}{1-e^{-\delta n_0}}
 \le2C_aQ(\delta n_0^2)
 \frac{E(\delta n_0^2/2)}{1-E(\delta n_0)}=B_n.
\]
Take $n_0=N+1\ge2$; then $\delta n_0\ge6/5$ and the latter denominator
is positive. Substituting $t=Tu$ in each retained-$n$ time tail gives
\[
 B_t=\sum_{n=1}^N2n^aT^3E(\delta n^2T)Q(\delta n^2T).
\]
The two tails are disjoint because the omitted-$n$ tail uses the whole
integral and the time tail uses only retained $n$.

For $j=0,\ldots,J-1$ and $\ell=0,\ldots,3$ take
\[
 m=2^j\left(1+\frac{2\ell+1}{8}\right),\qquad h=\frac{2^j}{8},
 \qquad T=2^J.
\]
These panels cover $[1,T]$ and have $h/m\le1/9$.
On $|z-m|=\rho=m/2$, principal log is analytic,
$\Re z\ge m/2$, and $|\arg z|\le1$.
For $f(z)=e^{-\pi n^2z/q}z^p$, the angular multiplier is at most
$e^{2/3}\le3$, whence
\[
 |f(z)|\le3B(m)E(\delta n^2m/2),\qquad
 B(m)=\max\{(3m/2)^2,(m/2)^{-2},1\}.
\]
Cauchy's coefficient estimate, cancellation of odd terms, and geometric
summation yield the $B_Q$ in Theorem~\ref{at:sample}.
Including zero-character or later discarded terms in that sum is harmless.
Gaussian, time, and geometric decay prove termination of the sequential
first-success searches with $B_n,B_t,B_Q\le\eta/8$.
For fixed q and $L=\max\{1,\log(1/\eta)\}$ their asymptotic scales can be
chosen as $N=O(\sqrt L)$, $T=O(L)$, $J=O(\log L)$, and $K=O(L)$.
Polynomial prefactors in $N,T,J$ contribute only $O(\log L)$ to the
logarithm of the error estimates; this is a truncation-size statement,
not a rational-arithmetic bit-complexity bound.

For real $t\in[m-h,m+h]$, $t\ge1$ and $t^{\Re p}\le(m+h)^2$.
The bound for a whole pair, including both powers, is therefore
$4hn^a(m+h)^2E(\delta n^2(m-h))$.
There are $P=4JN$ potential pairs; summing only bounds at most
$\eta/(8P)$ gives $B_{\rm discard}\le\eta/8$.
No partial recurrence or isolated exponential is discarded.
For active pairs retain the coefficient recurrence
\[
 r_{-1}=0,\quad r_0=1,\quad
 r_{j+1}=\frac{(h/m)(p-\lambda m-j)r_j-
                   \lambda m(h/m)^2r_{j-1}}{j+1},
 \qquad\lambda=\pi n^2/q.
\]
Coefficient comparison in $t f'=(-\lambda t+p)f$ proves that the
integrated Taylor polynomial is
\[
 2he^{p\log m-\lambda m}\sum_{j=0}^K\frac{r_{2j}}{2j+1}.
\]
Outward arithmetic encloses its exact value, including the signed
character weight and both powers.

Use an isolated p-bit dyadic context, with all primitive remainders
refined as p increases. Put $\delta_p=2^{-p-16}$.
Each arctangent series separately stops when its next-term bound is at
most $\delta_p$, giving total weighted analytic tail width at most
$20\delta_p$ in Machin's formula, before outward dyadic rounding.
Each reduced logarithm series separately has tail at most $\delta_p$;
for $x=2^kv$, the combined bound is $(1+|k|)\delta_p$, before rounding.
For the complex exponential, after reduction to $A\le1/2$, choose the
Taylor order so the bound
\[
 \frac{A^{H+1}}{(H+1)!}\frac1{1-A/(H+2)}.
\]
is at most $\delta_p$. Choose H by finite exact search; its denominator
is positive. Add a symmetric exponential error enclosure to each coordinate
before outward squaring. The fixed guard16 is not itself an acceptance
bound: the final complete rectangle radius must satisfy
$R_{\rm arith}\le\eta/4$. All interval caches must include p and preserve
exact input-type guards. The successful p and retry precisions/radii are
recorded. Individual series lengths are determined by these, the exact
inputs and the pinned source and are reconstructible by replay; they are
not explicitly stored.

For fixed N,J,K and exact target point, the expression has finitely many
operations and only nonzero fixed integer denominators. Each primitive
input is bounded for sufficiently large p; exponential halving counts
are consequently bounded. Exact primitive tails and dyadic rounding
widths tend to zero. Induction through the finite continuous expression,
including finitely many outward squarings, gives $R_{\rm arith}\to0$.
Doubling p after a failed comparison terminates. The argument also covers
true roots and nodes regenerated at every p, since their widths tend to zero.
A fixed nondegenerate s rectangle has no corresponding shrinking-input
claim.

\paragraph{A growing-order primitive estimate.}
To justify convergence despite a growing exponential Taylor order, write
$\epsilon_p=2^{-p}$, $\delta_p=2^{-p-16}$, and let $d(X)$ be the sum of a
rectangle's coordinate widths and $m(X)$ its $\ell^1$ magnitude bound.
The rounded reduced input $W$ satisfies $a=m(W)\le1/2+2\epsilon_p<3/5$
and $d(W)\le2^{-k}d(Z)+4\epsilon_p$; the remainder is at most $2^{-H}$,
so the first accepted order satisfies $H\le p+18$.
For the computed term boxes, outward multiplication and integer division give
$m_j\le a m_{j-1}+6\epsilon_p$ and
$d_j\le a d_{j-1}+m_{j-1}d(W)+12\epsilon_p$.
Thus $m_j\le a^j+15\epsilon_p$ and
$\sum_{j=0}^H d_j\le(25/4)d(W)+75H\epsilon_p$; adding the remainder
increases width by at most $4(\delta_p+\epsilon_p)$.
The pre-squaring width is therefore $O(d(Z)+p2^{-p})$ and its magnitude
is bounded (indeed below $3$ for $p\ge32$).
Each outward square has width at most $2m(X)d(X)+8\epsilon_p$;
the bounded number of squarings consequently preserves convergence to zero.
Finally, exact rational partial sums and endpoint rounding give Machin width
at most $20\delta_p+40\epsilon_p$, and logarithm width at most
$(1+|k|)(\delta_p+2\epsilon_p)$ for fixed $x=2^kv$.
These estimates justify the primitive-width step without assuming a fixed
Taylor order or asserting a global arithmetic complexity bound.


For serialization, initialize
\[
 B=\max\{32,\operatorname{bitlen}(\operatorname{den}\eta)
             -\operatorname{bitlen}(\operatorname{num}\eta)+8\}
\]
and increase B until $5\,2^{-B}\le\eta/8$.
Write $\widehat x=2^{-B}\lceil2^Bx\rceil$ for each of the five raw
components. Their exact reported sum is at most
$3\eta/4+5\,2^{-B}\le7\eta/8$, and every component remains an upper
bound. The output records B, the raw interval radius and complete rectangle,
N/J/K, discard threshold/pair counts, exact target metadata and source hashes.
These reconstruct all raw analytic bounds and pair sets without serializing
huge exact tail denominators. For the practical requests below B is87
and114, respectively. Exact endpoints remain rational; the mathematical
specification assumes an integer representation sufficient for the requested
size. Dynamic rounding removes the old fixed reporting floor. For rational s, validate its exact coordinates in closed D before
conversion: an outward dyadic enclosure of $\Im s=4/3$ extends beyond D,
but the analytic bounds concern the true point. Circle nodes satisfy the
strict exact bounds $-7/10\le\Re s\le17/10$, $|\Im s|\le6/5$.

The exact parameter-design checker gives
\[
\begin{array}{c|c|rrrr|r}
\eta&q&N&J&T&K&\text{active pairs}\\\hline
10^{-24}&3&12&7&128&19&71\\
10^{-24}&4&14&8&256&20&65\\
10^{-24}&5&16&8&256&20&111\\
10^{-32}&3&14&7&128&25&84\\
10^{-32}&4&16&8&256&26&72\\
10^{-32}&5&18&8&256&26&129
\end{array}
\]
with exact rational budgets and no special-function oracle.
The arithmetic precision is selected from the actual interval radius,
not inferred from this table. This extension retains the existing
classical theta identity; it adds no zero-location, positive determinant,
all-conductor, Lean or GRH claim. The following end-to-end adaptive readout supplies its own node-count,
alias, degree, Fourier and optimization choices.

% Generated shared reviewed material.
\section{An end-to-end adaptive analytic readout}\label{at:family}
The new driver changes the sample truncations, interval precision, circle
size and degree; it does not reuse the frozen labels to assert a smaller
error. Its scope is still the three stated primitive quadratic characters
and the fixed compact $|s-1/2|\le7/10$.

\begin{theorem}[Requested-error construction in the stated finite family]\label{at:requested}
Given a positive rational $\varepsilon$ and $q\in\{3,4,5\}$, the specified
exact-arithmetic algorithm terminates in the mathematical unbounded-resource
model and returns a normalized rational polynomial $P$ with
\[
 \sup_{|s-1/2|\le7/10}|P(s)-\Lambda_q(s)|\le\varepsilon.
\]
Its sample enclosures, rational optimizer updates and all error budgets are
explicit. This is an approximation theorem, with the stated classical theta
identity as an input; it supplies no zero-location or positive determinant
conclusion and is not a Lean theorem or a practical runtime bound.
\end{theorem}
\begin{proof}
Put $u=7/12$, $v=7/24$. Choose $d\ge0$ so
$10v^{d+1}/(1-v)\le\varepsilon/8$, and choose $m>d$ so
\[
 10\frac{2^{-m}}{1-2^{-m}}\frac1{1-v}\le\varepsilon/8.
\]
The driver rounds $m$ upward to a multiple of eight, which only decreases
this bound. Set the requested sample error
$\eta=\varepsilon(1-u)/8$. The adaptive sampler gives all $m$ true circle
nodes and labels with error at most $\eta$. Form their outward Fourier
rectangles through degree $d$, and refine the nodes and finite evaluation
arithmetic if their weighted midpoint radii exceed $\varepsilon/8$.
This refinement terminates: for these fixed finite $d,m$, node uncertainty
and every finite-expression arithmetic width tend to zero with precision.
Merely lifting fixed boxes would not suffice, so the driver regenerates the
nodes and samples on this retry.

Choose rational midpoint targets $t_n$ and execute the previously specified
quadratic-objective updates until
$\sum_{n=0}^d|b_n^k-t_n|_1u^n\le\varepsilon/4$.
Their geometric contraction proves finite termination. The oracle,
aliasing, Taylor and Fourier-midpoint terms are each at most
$\varepsilon/8$, and the optimization term is at most $\varepsilon/4$.
Hence the raw five-term sum is at most $3\varepsilon/4$.
Choose a serialization precision $B$ with $5\,2^{-B}\le\varepsilon/8$,
round each term upward, then add the five reported rationals exactly.
The printed bound is at most $7\varepsilon/8\le\varepsilon$.
The polynomial uses $((s-1/2)/(6/5))^n$, and the preceding Fourier/Cauchy
proof gives its asserted compact-uniform bound.
\end{proof}

The trace is losslessly encoded by the fixed target vector and a rational
scale extracted from each \emph{executed} coefficient state, together with
its objective, weighted residual and full-state hash. Every coefficient
can be recovered exactly; a separate replay expands the states and checks
every update. The producer does not generate its trace from a closed form.
The saved objective remains the rational Fourier-midpoint quadratic, with
its perturbation from ideal-root least squares explicitly budgeted.

% Generated shared reviewed material.
\section{Termination, compression and the boundary of the guarantee}\label{at:family-details}
The degree and circle-size choices use only rational comparisons. Since
$0<v<1$, the Taylor test succeeds at finite $d$; since $2^{-m}\to0$, the
alias test succeeds at finite $m>d$. The common growth bound10 is safe for
all three families; the certificate also recomputes the previously proved
exact conductor-specific majorants below $4,10,6$.

At fixed $d,m,\eta$, the analytic truncation parameters are fixed while
precision is doubled. The computed sample centers stay bounded: their
underlying finite expression is fixed and the true completed-function
values are bounded on the sample circle. Its arithmetic widths tend to
zero, as do the regenerated true-root enclosures. Fourier computation is
a finite composition of additions, products and division by the positive
integer $m$. Its midpoint-radius sum therefore tends to zero. This proves
the driver's additional precision retry, without pretending that changing
a box's nominal precision shrinks its uncertainty.

The accepted sampler error may remain above its actual arithmetic width
because its analytic tails and output bounds are fixed for this $\eta$.
That causes no problem: the oracle allocation uses the total certified
sample radius, whereas the Fourier-midpoint allocation uses only finite
coefficient-rectangle width. No analytic sample error is silently moved
into the latter term or removed by numerical convergence.

For the stored training trace, the exact loop gives $b^{k+1}=3b^k/4+t/4$
from $b^0=0$. Inductively every state is a rational scalar multiple of the
same target vector. The producer first executes the actual vector updates,
then extracts this scalar from a nonzero target coordinate, checks all
coordinates, and records a canonical full-vector SHA-256 hash. If the
entire target vector is zero, its zero state needs only scale0. The compressed
record is thus an exact representation of the executed states, not a
floating-point sketch. Hashes are integrity aids; exact vector reconstruction
and the independently checked recurrence establish the mathematical claim.

Coefficient export uses exact rational numerators and denominators. A
chunked decimal integer encoder/parser handles endpoints beyond Python's
usual digit-conversion guard without changing global interpreter settings.
No fixed $180$-bit serialization floor remains. The decimal graphics are
separate approximate projections, not polynomial coefficients to which the
exact certificate may be applied without a conversion-error bound.

All producing source bytes are captured, copied into fresh output paths and
checked unchanged. Historical fixed producers are pinned and untouched.
For the new modules the run also checks the import-time source snapshots,
so a stale imported module cannot be attributed to later source bytes.
After output creation, a failed evaluation, final provenance check or
serialization emits uncertified failure metadata. Invalid requests and stale
import-time sources are refused before creating output artifacts. Old
outputs are never overwritten. Extreme requests can be
computationally impractical: the termination proof gives no machine-time or
memory guarantee, and finite test runs establish only their reported cases.

Nothing requires the output polynomials to be real-rooted, to equal an
arithmetic Fredholm determinant, or to have a self-adjoint realization.
Consequently arbitrary compact approximation of these known functions does
not imply GRH. Uniformity in changing conductor, general Hecke/Dedekind/
elliptic/Artin scopes, the independent determinant comparison, and actual
analytic/category Lean objects remain separate obligations.


\chapter{The original-family Hecke moment and exact masked transfer}
% Generated shared reviewed material.
% Accepted prose AFM proof; explicit external inputs, not a Lean theorem.
\section{An arithmetic-family moment with every original row retained}
\label{mm:core}
Let $K=\mathbb Q(\sqrt{-3})$, $\mathcal O=\mathbb Z[\omega]$, and write
$N$ for the ideal or element norm. Fix a finite-order Hecke character
$\nu$ and a finite set $S$ of prime ideals containing those above $6$ and
those dividing its conductor. For $0\ne u\in\mathcal O$, let $\kappa_u$
be the character of $K(u^{1/6})/K$ defined using arithmetic Frobenius, and
put $\psi_u=\operatorname{prim}(\nu\kappa_u)$. Here
$\operatorname{prim}$ means the underlying primitive Hecke character:
it is formed before introducing any divisor variable or evaluating
conductor-dependent zero extensions. The rows below are all actual
Eisenstein elements, including units and prime powers supported in $S$.

The external inputs are abelian Kummer reciprocity and character
uniqueness, the primitive Hecke functional equation and its usual
vertical growth, Stirling's formula, scalar Dedekind-zeta convexity, and
Blomer--Goldmakher--Louvel's large sieve for their specified sixth-order
family. We use the pinned
\href{https://arxiv.org/pdf/1112.1650v1}{arXiv:1112.1650v1},
Theorem~1.3, Lemma~2.1 and Appendix~II; its Section~5 supplies the
all-ideal decomposition whose fixed-twist proof is redone here.
No statement about all primitive characters of comparable conductor,
Artin holomorphy, Lindel\"of bounds or zero locations is assumed.

\begin{theorem}[The original-family numerator moment]
\label{mm:afm}
For every $\varepsilon>0$, uniformly for $H\ge1$ and $t\in\mathbb R$,
\begin{equation}\label{mm:central}
 \sum_{0<Nu\le H}|L(1/2+it,\psi_u)|^2
 \ll_{\nu,S,\varepsilon}H^{1+\varepsilon}(1+|t|)^{1+\varepsilon}.
\end{equation}
For every fixed $0<\delta<1/2$, the same estimate holds with
$1/2+it$ replaced by $1/2+\delta+it$, with the constant also depending
on $\delta$. If the central bound has constant $C_0$ at an exponent
$0<\varepsilon_0\le1/4$, an admissible off-central constant is
\begin{equation}\label{mm:poisson-constant}
 200 C_0\left(\frac{3/2+\delta}{1/2-\delta}\right)^2.
\end{equation}
\end{theorem}
\begin{proof}
For the canonical BGL character $\vartheta_{\mathfrak a}$,
Appendix~II gives a radicand $\alpha_{\mathfrak a}$ with
$(\alpha_{\mathfrak a})_S=\mathfrak a\mathfrak g^{-6}$.
Since $\mathcal O_S$ is principal, comparison with any literal $u$
whose outside-$S$ ideal is $\mathfrak a$ gives
$\kappa_u=\eta\vartheta_{\mathfrak a}$ for one of finitely many fixed
$S$-unit Kummer twists $\eta$. This is an equality of primitive
characters, established outside a finite set and extended by uniqueness.
For each outside-$S$ ideal $\mathfrak a$ and $S$-supported integral
ideal $\mathfrak b$, at most six elements generate
$\mathfrak a\mathfrak b$. The resulting $S$-supported multiplicity sum
converges at exponent $1+\varepsilon$.

For each fixed $\xi=\nu\eta$, split the canonical parameter as
$\mathfrak a=\mathfrak a_0\mathfrak a_1\mathfrak a_2^6$, where
$\mathfrak a_0$ is squarefree, $\mathfrak a_1$ is squarefull and
sixth-power-free, and $(\mathfrak a_0,\mathfrak a_1)=1$.
The primitive conductor is at most a fixed multiple of
$N(\mathfrak a_0\mathfrak a_1)$.
The common-gamma approximate functional equation, with
$X=H/(N\mathfrak a_2)^6$ and $T=1+|t|$, has common maximal length
$B=X^{1/2}T(HT)^\eta$ for any sufficiently small fixed $\eta>0$.
Mellin inversion removes its row-dependent conductor weight.
Split a column as $\mathfrak b_S\mathfrak b_1\mathfrak b_2$, with
$\mathfrak b_S$ supported in $S$, $\mathfrak b_1$ squarefree and
$\mathfrak b_2$ squarefull, the last two coprime and prime to $S$.
The original primitive values at $\mathfrak b_S$ remain bounded row
multipliers. The $S$-supported and squarefull column weights are summable
at $1/2+\eta$. Multiplying squarefree column coefficients by the fixed
$\xi$ preserves their squared norm, whatever its finite order.
BGL's large sieve therefore bounds the energy for fixed
$\mathfrak a_1,\mathfrak a_2$ by
\[
 (HT)^{O(\eta)}\{M+B+(MB)^{2/3}\},\qquad
 M=X/N\mathfrak a_1.
\]
Squarefull counting and summation give
$X+B\sqrt X+(XB)^{2/3}\ll (HT)^{O(\eta)}XT$.
The $\mathfrak a_2$ sum converges at exponent $6$.
Principal rows are separated before the pole-free functional equation
and bounded by scalar zeta convexity. This proves~\eqref{mm:central}.

For the off-central statement divide each $L(s,\psi)$ by
$(s+1)^{5/4}$ and, on principal rows, also multiply by $(s-1)/(s+1)$.
These are bounded analytic functions on $\operatorname{Re}s>1/2$,
continuous on the boundary. Poisson representation and Jensen's
inequality transfer the central squared moment. The convolution of
$P_\delta(y)=\delta/[\pi(\delta^2+y^2)]$ with
$(1+|t|)^{-3/2+\varepsilon_0}$ is at most eleven times that weight
when $0<\varepsilon_0\le1/4$. Restoring the denominator and the
principal pole-removal factor gives~\eqref{mm:poisson-constant}.
The extended edition gives the detailed derivation with all primitive local values.
\end{proof}

\paragraph{Separate fixed-unitary extension.}
The moment theorem also holds for any one fixed unitary Hecke character
$\nu$. For precision choose
$\nu(\mathfrak a)=\nu_0(\mathfrak a)(N\mathfrak a)^{-i\tau}$, with
angular parameter $m\in\mathbb Z$ and completed gamma factor
\[
 \gamma_\nu(s)=2(2\pi)^{-(s+|m|/2+i\tau)}
                    \Gamma(s+|m|/2+i\tau).
\]
All finite Kummer twists have this same infinity type. Use
$T_\nu=1+|t+\tau|+|m|$, comparable to $1+|t|$ with
$\nu$-dependent constants. Pure norm rows have
$L(s,\psi)=\zeta_K(s+i\tau)$ and are separated before the pole-free
functional equation; their pole at $1-i\tau$ is removed by
$(s-1+i\tau)/(s+1+i\tau)$. No uniformity in $m$, $\tau$ or the
conductor of $\nu$ is asserted. This extension of the moment is separate
from the following finite-order type-I residue formula.

\subsection{The exact masked remainder and the remaining loss}
Retain the original zero-extended symbol $\chi_d(u)=(u/d)_6$ on good
squarefree ideals. Let $\mu$ be the ideal M\"obius function and set
\[
 \beta_U(d)=\sum_{ab=d\atop Na,Nb\le U}\mu(a)\mu(b),\quad
 T(u,d)=S\cup\operatorname{supp}(u)\cup\operatorname{supp}(d).
\]
Let $\psi_{u,2}=\operatorname{prim}(\psi_u^2)$, formed independently
of squaring any zero sentinel. With $\mathfrak f_u$ and
$\mathfrak f_{u,2}$ the primitive conductors, the full-mask series is
\begin{align}
 F_{u,d}(s)
 &=\frac{L(s,\psi_u)}{L(2s,\psi_{u,2})}C_{T(u,d)}(s),\label{mm:masked}\\
 C_T(s)
 &=\frac{\prod_{p\in T,\ p\nmid\mathfrak f_u}
              (1-\psi_u(p)(Np)^{-s})}
         {\prod_{p\in T,\ p\nmid\mathfrak f_{u,2}}
              (1-\psi_{u,2}(p)(Np)^{-2s})}.\notag
\end{align}
Initially it equals
$\sum_{c\ \mathrm{squarefree},\ (c,T)=1}\psi_u(c)(Nc)^{-s}$
for $\operatorname{Re}s>1$. Every prime of $T$ remains deleted,
including those forgotten by either primitive conductor.

For finite-order $\nu$, $D\ge1$, $U\ge1$, $U^2\le D$, and
$W\in C_c^\infty((0,\infty))$ supported in $[1,2]$, put
\[
 \widehat W(s)=\int_0^\infty W(x)x^{s-1}\,dx,\qquad
 \mathcal M_4(W)=\sum_{j=0}^4\sup_x|(x\partial_x)^jW(x)|.
\]
On $\sigma=1/2+\delta$, the exact residue-subtracted type-I remainder is
\begin{equation}\label{mm:remainder}
 E_u=\frac{D^\delta}{2\pi}\int_{\mathbb R}
 \widehat W(\sigma+it)D^{it}
 \sum_{d\ \mathrm{squarefree}\atop(d,S)=1}
 \frac{\beta_U(d)\nu(d)\chi_d(u)}{(Nd)^{\sigma+it}}
 F_{u,d}(\sigma+it)\,dt.
\end{equation}
This definition retains the actual $\beta_U$, all phases, $d$-masks,
conductors and the independently primitive square.

\begin{theorem}[An improved masked residual estimate]
\label{mm:residual}
Under the preceding finite-order hypotheses, for every $\varepsilon>0$,
\begin{equation}\label{mm:energy}
 \sum_{0<Nu\le H}|E_u|^2
 \ll_{\nu,S,\delta,\varepsilon}
 \mathcal M_4(W)^2D^{2\delta}U^{2-4\delta}H(HD)^\varepsilon.
\end{equation}
\end{theorem}
\begin{proof}
The same deleted set in both factors of~\eqref{mm:masked} gives
$|F_{u,d}(\sigma+it)|\ll (Nu\,Nd)^\eta|L(\sigma+it,\psi_u)|$:
the reciprocal incomplete Euler product at $2\sigma>1$ is bounded
by $\zeta_K(2\sigma)$, and its remaining finite numerator is bounded
by a small power of its radical. Apply Minkowski to~\eqref{mm:remainder}
in the exact row space and use Theorem~\ref{mm:afm}. The actual coefficient
sum is at most
\[
 \sum_d\frac{|\beta_U(d)|}{(Nd)^{\sigma-\eta}}
 \le\left(\sum_{Na\le U}(Na)^{-\sigma+\eta}\right)^2
 \ll U^{1-2\delta+2\eta}.
\]
Four Mellin integrations by parts bound the remaining height integral
by a constant times $\mathcal M_4(W)$. Square the result and absorb
the small losses using $U^2\le D$.
\end{proof}

At $H=D^{3/4}$ and $U=D^a$ with fixed $0<a\le1/2$, the upper bound's excess
over $H$ is $2\delta+(2-4\delta)a$, up to small losses; as
$\delta\downarrow0$ it still leaves $2a$. This is a limitation of the
upper bound, not an energy lower bound. The extracted principal residue,
the original short term at smaller scales and the signed type-II minus
type-I covariance are not controlled by~\eqref{mm:energy}. The next
arithmetic task is to estimate the complete $d$-weighted moment before
its absolute coefficient sum, or establish sufficient signed covariance.
No positive arithmetic determinant, zero-location theorem or GRH proof
follows. The analytic argument uses the external classical inputs above;
finite diagnostics and selected existing Lean local identities do not
formalize the Hecke functions, approximate functional equation, Poisson
transfer or these moment bounds.

% Generated shared reviewed material.
% Full derivation of the scoped arithmetic-family moment and masked remainder.
\section{Family identification, the approximate functional equation and masks}
\label{mm:details}
We prove the analytic statements of Section~\ref{mm:core} from the
explicit classical inputs stated there. Write $I(S)$ for the nonzero
integral ideals prime to $S$, and $\mathcal O_S$ for the ring of
$S$-integers. Constants may depend on the fixed field, $S$, $\nu$ and
any declared positive small exponent. No such constant is asserted
uniform in the conductor or infinity type of a varying $\nu$.

\subsection{A finite comparison with the literal Kummer rows}
Use Appendix~II of BGL, with $n=6$. Its fixed local modulus and
representatives may be chosen once, since $S$ contains the primes above
$6$ and $\mathcal O_S$ has class number one. For each
$\mathfrak a\in I(S)$ the construction writes
\[
 \mathfrak a=(x)E\mathfrak g^6,\qquad
 \alpha_{\mathfrak a}=xm_E,\qquad
 (m_E)_S=E,
\]
so $(\alpha_{\mathfrak a})_S=\mathfrak a\mathfrak g^{-6}$.
Choose $y\in K^\times$ generating $\mathfrak g$ as an
$\mathcal O_S$-ideal. If $(u)_S=\mathfrak a$, then
\[
 e=\alpha_{\mathfrak a}y^6/u\in\mathcal O_S^\times,\qquad
 [\alpha_{\mathfrak a}]=[eu]\quad\hbox{in }K^\times/K^{\times6}.
\]
There is a precision here: Appendix~II allows choosing $\mathfrak g$
relative to a tested denominator. Fix one admissible global decomposition
first. Its radicand gives the canonical character
$\vartheta_{\mathfrak a}$ on all primes outside the finite set
$S\cup\operatorname{supp}(\mathfrak a\mathfrak g)$ where that
choice applies. Uniqueness of Hecke characters, equivalently Chebotarev
in the finite abelian quotient, then identifies the underlying primitive
characters everywhere. With matching arithmetic-Frobenius conventions,
\begin{equation}\label{mm:unit-comparison}
 \kappa_u=\kappa_{e^{-1}}\vartheta_{\mathfrak a}.
\end{equation}
If the residue-symbol convention is conjugated, conjugate the canonical
family and its coefficient sequence; the same large sieve follows.

The group $\mathcal O_S^\times/(\mathcal O_S^\times)^6$ is finite.
Thus~\eqref{mm:unit-comparison} uses one of a finite set of fixed
Kummer twists $\eta$, all unramified outside $S$ and of bounded
$S$-conductor. This is a comparison of primitive characters, not a
pointwise multiplication rule for independently zero-extended arrays at
bad primes. In particular, a conductor drop must not be replaced by a
product of artificial zeros.

Since $\mathcal O$ is principal with six units, for fixed
$\mathfrak a\in I(S)$ and an $S$-supported integral ideal
$\mathfrak b$ there are at most six $u$ satisfying
$(u)=\mathfrak a\mathfrak b$. For any nonnegative row quantity, the
literal moment is consequently at most
\begin{equation}\label{mm:literal-sum}
 6\sum_{\mathfrak b\ S\text{-supported}\atop N\mathfrak b\le H}
   \sum_{\eta}\sum_{\mathfrak a\in I(S)\atop N\mathfrak a\le H/N\mathfrak b}
 \left|L\bigl(s,\operatorname{prim}(\nu\eta
                       \vartheta_{\mathfrak a})\bigr)\right|^2.
\end{equation}
An ideal bound $X^{1+\varepsilon}T^{1+\varepsilon}$ therefore transfers
to all literal rows, since
\[
 \sum_{\mathfrak b\ S\text{-supported}}(N\mathfrak b)^{-1-\varepsilon}
 =\prod_{p\in S}(1-(Np)^{-1-\varepsilon})^{-1}<\infty.
\]
No quotient of the original unit rows or deletion of $S$-prime powers
has occurred.

\subsection{Fixed twists, exact tame conductors and the squarefree sieve}
Fix $\xi=\nu\eta$ and write
$\rho_{\mathfrak a}=\operatorname{prim}
(\xi\vartheta_{\mathfrak a})$. Its fixed conductor is supported in
$S$ before the Kummer twist. On columns prime to $S$, $|\xi(\mathfrak b)|=1$.
BGL Theorem~1.3 thus gives, for any fixed $\mathfrak a_1$,
\begin{align}\label{mm:ls}
 &\sum_{\mathfrak a_0\in I(S)\ \mathrm{squarefree}\atop N\mathfrak a_0\le M}
 \left|\sum_{\mathfrak b_1\in I(S)\ \mathrm{squarefree}\atop N\mathfrak b_1\le B}
 \lambda_{\mathfrak b_1}\xi(\mathfrak b_1)
 \vartheta_{\mathfrak a_1}(\mathfrak b_1)
 \vartheta_{\mathfrak a_0}(\mathfrak b_1)\right|^2\notag\\
 &\hspace{1cm}\ll (MB)^\varepsilon
 \{M+B+(MB)^{2/3}\}\sum_{\mathfrak b_1}|\lambda_{\mathfrak b_1}|^2.
\end{align}
The first fixed multiplier preserves the coefficient norm; the second
has modulus at most one. Restricting the positive outer row sum is
harmless. A coprimality condition with another fixed ideal belongs to
$\lambda$. In particular, $\nu$ need not have order dividing six and
is not asserted to belong to BGL's finite ray-class character set.

There is a unique valuation decomposition
\[
 \mathfrak a=\mathfrak a_0\mathfrak a_1\mathfrak a_2^6,
 \quad \mathfrak a_0\ \hbox{squarefree},\quad
 \mathfrak a_1\ \hbox{squarefull and sixth-power-free},\quad
 (\mathfrak a_0,\mathfrak a_1)=1.
\]
Indeed, for $v_p(\mathfrak a)=6q+r$, $0\le r\le5$, put the residual
prime in $\mathfrak a_0$ if $r=1$, and its exponent $r$ in
$\mathfrak a_1$ if $2\le r\le5$. BGL Lemma~2.1 gives
$\vartheta_{\mathfrak a\mathfrak b^6}=\vartheta_{\mathfrak a}$ and
$\vartheta_{\mathfrak a_0\mathfrak a_1}
=\vartheta_{\mathfrak a_0}\vartheta_{\mathfrak a_1}$ on ideals
prime to $S$ in this sixth-power-free case. Those are the identities
needed here; no arbitrary multiplicativity assertion is substituted.

At a prime outside $S$, tame Kummer inertia at residual exponent $r$
has order $6/\gcd(6,r)$. It is nontrivial for $r=1,\ldots,5$ and
has conductor exponent one. The fixed $\xi$ is unramified there, so
cannot cancel that inertia. At primes in $S$ the local sixth-power
classes are finite, hence only finitely many primitive conductor types
occur after multiplication by the fixed $\xi$. Consequently
\begin{equation}\label{mm:conductor}
 Q_{\rho_{\mathfrak a}}=|D_K|N\mathfrak f_{\rho_{\mathfrak a}}
 \le C_{\xi,S}N(\mathfrak a_0\mathfrak a_1),
\end{equation}
independently of $\mathfrak a_2$. The equality of the conductor's
outside-$S$ part with $\operatorname{rad}(\mathfrak a_0\mathfrak a_1)$
does not identify that conductor with the original larger deleted set.

\subsection{Removing the row-dependent functional-equation weights}
First take finite-order $\nu$. All characters then have trivial infinity
type at the complex place. The primitive completion may be normalized as
\[
 \Lambda(s,\rho)=Q_\rho^{s/2}\gamma_K(s)L(s,\rho),\qquad
 \gamma_K(s)=2(2\pi)^{-s}\Gamma(s),\qquad
 \Lambda(s,\rho)=\epsilon_\rho\Lambda(1-s,\bar\rho),
 \quad |\epsilon_\rho|=1.
\]
For nonprincipal rows the completed function is entire. Apply this
functional equation to the Mellin integral with even auxiliary factor
$G(z)=e^{z^2}$ and denominator $z$, shifting across $z=0$. At
$s=1/2+it$ this gives two weighted Dirichlet series, truncated below to
polynomials; the dual prefactor has modulus one. Principal rows are excluded from
this pole-free step and treated separately below.

For fixed $\mathfrak a_2$, let $X=H/(N\mathfrak a_2)^6\ge1$ and
$T=1+|t|$. An explicit normalized kernel is
\begin{equation}\label{mm:afe-kernel}
 V_t(x)=\frac1{2\pi i}\int_{(\eta)}
 e^{z^2}\frac{\Gamma(s+z)}{\Gamma(s)}T^{-z}x^{-z}\,\frac{dz}{z},
 \qquad x=\frac{2\pi N\mathfrak b}{\sqrt{Q_\rho}T},
\end{equation}
with $\eta>0$ small. Its factors combine to
$(\sqrt{Q_\rho}/(2\pi N\mathfrak b))^z$ in the actual completed-L
kernel; there is no extra factor of $T$. The denominator $z$ supplies
the residue at zero.
On $z=\eta+iy$, Stirling and $|t+y|\le T(1+|y|)$ give
\[
 \left|\frac{\Gamma(1/2+\eta+i(t+y))}
                 {\Gamma(1/2+it)T^{\eta+iy}}\right|
 \ll_\eta (1+|y|)^{C_\eta}e^{\pi|y|/2}.
\]
The factor $e^{-y^2}$ makes this an integrable majorant uniformly in
$t$. The analogous bound holds on every fixed positive shifted line
and for the dual gamma kernel. Moving far enough to the right gives
arbitrary power decay beyond a common total-ideal cutoff
\begin{equation}\label{mm:common-cutoff}
 B'=X^{1/2}T(HT)^\eta,
\end{equation}
up to fixed $\xi,S,K$ constants. For any requested $A>0$, a sufficiently
large positive kernel shift, depending also on $\eta$, makes the
truncation error $O((HT)^{-A})$ uniformly in all rows in this block.
Ideal counting $O_K(Y)$ bounds the absolute tails. Such errors are
negligible after the finite row sums by taking $A$ larger if needed.

Insert~\eqref{mm:afe-kernel} on $\operatorname{Re}z=\eta$.
The factor $(\sqrt{Q_\rho}T)^z$ is independent of the column and is a
row multiplier of modulus at most $(CXT^2)^{\eta/2}$.
Its imaginary part is a phase. The remaining column coefficient is
$(N\mathfrak b)^{-1/2-\eta-it-iy}$, apart from fixed constants.
Minkowski or Cauchy with the integrable gamma majorant therefore removes
the row-dependent weight at cost $(HT)^{O(\eta)}$ in the squared norm.
No independence of the root numbers is assumed or needed.

For each column use the unique decomposition
\[
 \mathfrak b=\mathfrak b_S\mathfrak b_1\mathfrak b_2,
 \quad \mathfrak b_S\ S\hbox{-supported},\quad
 \mathfrak b_1\ \hbox{squarefree},\quad
 \mathfrak b_2\ \hbox{squarefull},\quad
 (\mathfrak b_1,\mathfrak b_2)=1,
\]
where $\mathfrak b_1,\mathfrak b_2\in I(S)$; a prime of outside-$S$
valuation one goes in $\mathfrak b_1$, and a prime of valuation at least
two goes with its full exponent in $\mathfrak b_2$. For fixed
$\mathfrak b_S,\mathfrak b_2$ retain
$\rho_{\mathfrak a}(\mathfrak b_S)$ as the original primitive row
value. On the good factors the exact canonical identity applies.
The remaining factors involving $\mathfrak b_2$, including
$\vartheta_{\mathfrak a_0}(\mathfrak b_2)$, are bounded row
multipliers. Their moduli are at most one; only then are they discarded
in an inequality. In particular the value at $\mathfrak b_S$ is not
replaced by the product of zero-extended primitive arrays.

For the $\mathfrak b_1$ polynomial the coefficient sequence is
independent of $\mathfrak a_0$; it includes the indicator
$(\mathfrak b_1,\mathfrak b_2)=1$ and the common cutoff
$N\mathfrak b_1\le B'/N(\mathfrak b_S\mathfrak b_2)$.
Padding it by zeros up to $B'$ is legitimate. Its squared norm is at most
$\zeta_K(1+2\eta)$. Minkowski over $\mathfrak b_S,\mathfrak b_2$
is legitimate because
\[
 \sum_{\mathfrak b_S,\mathfrak b_2}
       N(\mathfrak b_S\mathfrak b_2)^{-1/2-\eta}<\infty.
\]
The finite Euler product handles $\mathfrak b_S$; the squarefull series
converges above $1/2$. Apply~\eqref{mm:ls} to each such polynomial and
to the conjugate dual family. For fixed $\mathfrak a_1,\mathfrak a_2$
the resulting energy is
\begin{equation}\label{mm:block-energy}
 (HT)^{O(\eta)}\{M+B'+(MB')^{2/3}\},\qquad
 M=X/N\mathfrak a_1.
\end{equation}
The additional condition $(\mathfrak a_0,\mathfrak a_1)=1$ simply
restricts the positive row sum. The large-sieve small exponent can be
chosen sufficiently small along with $\eta$ and absorbed in the
indicated small loss.

\subsection{All-ideal summation and principal rows}
Every squarefull ideal $\mathfrak c$ has a unique expression
$\mathfrak d^2\mathfrak e^3$ with $\mathfrak e$ squarefree.
This follows primewise by using exponent zero or one in $\mathfrak e$
according to the parity of an exponent at least two. Therefore
\[
 \#\{\mathfrak c\ \mathrm{squarefull}:N\mathfrak c\le X\}
 \ll_K X^{1/2},\qquad
 \sum_{\mathfrak c\ \mathrm{squarefull}}(N\mathfrak c)^{-s}<\infty
 \quad(s>1/2).
\]
For the count, sum
$O((X/(N\mathfrak e)^3)^{1/2})$ and use
$\zeta_K(3/2)<\infty$. Summing~\eqref{mm:block-energy} over
$\mathfrak a_1$ gives
\[
 (HT)^{O(\eta)}
 \{X+B'\sqrt X+(XB')^{2/3}\}
 \ll (HT)^{O(\eta)}XT,
\]
since the squarefull series converges at $1$ and $2/3$,
$B'=\sqrt X\,T(HT)^\eta$, and $T^{2/3}\le T$ for $T\ge1$.
The $\mathfrak a_2$ sum is bounded by
$H\sum_{\mathfrak a_2}(N\mathfrak a_2)^{-6}=H\zeta_K(6)$.
Choosing all internal losses small for the requested final
$\varepsilon$ proves the fixed-twist all-ideal central moment.

There are $O_K(H)$ ideal rows in total. On any row where the primitive
character is principal its $L$-function is $\zeta_K$, and ordinary
scalar convexity gives
$|\zeta_K(1/2+it)|^2\ll_\varepsilon T^{1+\varepsilon}$.
This treats those rows without a pole-free approximate functional
equation. Finally~\eqref{mm:literal-sum} proves~\eqref{mm:central}
for every actual nonzero Eisenstein element.

\subsection{Poisson transfer with an explicit constant}
For a nonprincipal row define
\[
 f_\psi(s)=\frac{L(s,\psi)}{(s+1)^{5/4}},
\]
using the holomorphic logarithm on $\operatorname{Re}(s+1)>0$.
On a principal row additionally multiply by $(s-1)/(s+1)$.
Analytic continuation and ordinary degree-two convexity, used separately
at each fixed conductor, show that $f_\psi$ is bounded analytic on the
right half-plane $\operatorname{Re}s>1/2$, with continuous boundary
and decay at infinity. The principal pole is removed before applying
Poisson representation to its real and imaginary parts. Jensen gives
\[
 |f_\psi(1/2+\delta+it)|^2
 \le\int_{\mathbb R}P_\delta(t-v)
             |f_\psi(1/2+iv)|^2\,dv.
\]
The exact row family is finite for fixed $H$, so its sum interchanges
with the integral without a limiting-family issue. The boundary
pole-removal factor has modulus at most one.

Take $0<\varepsilon_0\le1/4$ and
$\alpha=3/2-\varepsilon_0\in[5/4,3/2)$. Since
$|3/2+iv|\ge(1+|v|)/\sqrt2$, the central squared moment of the
$f_\psi$ is bounded by
$3C_0H^{1+\varepsilon_0}w_\alpha(v)$, where
$w_\alpha(v)=(1+|v|)^{-\alpha}$.
For $|t|\ge2$, split $|v-t|\le|t|/2$ from its complement. The
near contribution to $P_\delta*w_\alpha$ is at most
$2^\alpha w_\alpha(t)<3w_\alpha(t)$. In the far region use
$P_\delta(t-v)\le4\delta/(\pi t^2)$ and
$\int w_\alpha=2/(\alpha-1)$, giving at most
\[
 \frac{2\delta3^\alpha}{\pi(\alpha-1)}w_\alpha(t)
 \le8w_\alpha(t).
\]
Here $(1+|t|)^\alpha/t^2\le3^\alpha/4$ for $|t|\ge2$.
For $|t|\le2$, the convolution is at most one and hence at most
$3^\alpha w_\alpha(t)<6w_\alpha(t)$. Thus
$P_\delta*w_\alpha\le11w_\alpha$ uniformly for $0<\delta<1/2$.
Restoring the denominator costs less than
$6(1+|t|)^{5/2}$, and the inverse principal factor is bounded by
$(3/2+\delta)/(1/2-\delta)$. The product $11\cdot3\cdot6=198<200$
proves~\eqref{mm:poisson-constant}. The constant is uniform as
$\delta\downarrow0$, but not as the contour approaches the principal
pole line $\operatorname{Re}s=1$. No zero-free hypothesis appears.

\subsection{The separately scoped unitary infinity type}
Let $\nu(\mathfrak a)=\nu_0(\mathfrak a)(N\mathfrak a)^{-i\tau}$
be fixed and unitary, with angular parameter $m$ and gamma factor
$\gamma_\nu$ stated in Section~\ref{mm:core}. The external primitive
functional equation is explicitly
\[
 Q_\rho^{s/2}\gamma_\nu(s)L(s,\rho)
 =\epsilon_\rho Q_\rho^{(1-s)/2}
                 \gamma_{\bar\nu}(1-s)L(1-s,\bar\rho),
 \qquad |\epsilon_\rho|=1.
\]
A constant norm-twist phase in another completion convention is absorbed
in $\epsilon_\rho$. All finite Kummer twists share this infinity type.
Replace $T$ in the kernel by $T_\nu=1+|t+\tau|+|m|$.
The numerator gamma ratio now has real base $1/2+|m|/2>0$ and
vertical parameter $t+\tau$; its conjugate dual has parameter
$-(t+\tau)$. Stirling gives the same integrable Gaussian-majorized
kernel on each fixed shifted line. Constants can depend on $m,\tau$;
$T_\nu$ is comparable to $T$ for this one fixed $\nu$.
The finite-prime conductor analysis, modulus-one coefficient twist,
all-ideal decomposition and all original primitive row values remain
unchanged.

A pure norm row has $L(s,\rho)=\zeta_K(s+i\tau)$; exclude it from
the pole-free approximate functional equation and use scalar convexity
at $t+\tau$. In the Poisson step cancel its actual pole by
\[
 h_\tau(s)=\frac{s-1+i\tau}{s+1+i\tau}.
\]
This factor is analytic in the relevant right half-plane. On its
boundary $|h_\tau|\le1$, while inside
$|h_\tau^{-1}|\le(3/2+\delta)/(1/2-\delta)$.
Hence the same transfer factor $200$ applies to the now
$\nu$-dependent central constant. If the opposite norm convention is
chosen, replace $\tau$ by $-\tau$ in the gamma factor and pole alike.
This proves the separate fixed-unitary moment assertion; it does not
claim uniformity in a family of infinity types or replace the
finite-order principal-residue expression below.

\subsection{Keeping the full type-I mask during transfer}
For $\sigma=1/2+\delta$, define $L^T$ by removing every prime of
$T(u,d)$, not just conductor primes. The identity~\eqref{mm:masked}
is equivalently
\[
 F_{u,d}(s)=\frac{L(s,\psi_u)}{L^{T(u,d)}(2s,\psi_{u,2})}
 \prod_{p\in T(u,d),\ p\nmid\mathfrak f_u}
                         (1-\psi_u(p)(Np)^{-s}).
\]
Absolute convergence at $2\sigma>1$ bounds the reciprocal incomplete
product by $\zeta_K(2\sigma)$, without requiring any zero-free theorem
in the critical strip. For every fixed $\eta>0$,
\[
 \prod_{p\in T}(1+(Np)^{-\sigma})
 \ll_{\sigma,\eta}(N\operatorname{rad}T)^\eta,
 \qquad N\operatorname{rad}T\le C_S Nu\,Nd.
\]
Indeed $1+q^{-\sigma}\le q^\eta$ beyond a fixed threshold, and only
finitely many prime ideals lie below it. Therefore
\begin{equation}\label{mm:mask-bound}
 |F_{u,d}(\sigma+it)|
 \ll_{\nu,S,\delta,\eta}(Nu\,Nd)^\eta
                              |L(\sigma+it,\psi_u)|.
\end{equation}
This is a bounded inequality after the exact mask identity, not an
identity that removes original zeros. For example, with principal $\nu$
and $u=\pi^3$ or $\pi^6$ at $N\pi=7$, the primitive local ratio,
restoration and full masked factor at $s=1$ are respectively
\[
 \begin{array}{c|ccc}
  v_\pi(u)&\hbox{primitive ratio}&\hbox{restoration}&\hbox{masked factor}\\\hline
 3&48/49&49/48&1\\
 6&8/7&7/8&1
 \end{array}
\]
The squared primitive character in the first row is unramified at
$\pi$ although the original primitive character is ramified there.
Squaring a zero sentinel would give the wrong denominator.

For finite-order $\nu$, the normalized type-I piece is the retained
$d$ sum of the smoothed squarefree series at $Y=D/Nd$, divided by
$\sqrt D$. Its possible pole is only at $s=1$, on principal rows.
Contour shifting with the full mask gives $I_u=P_u+E_u$, with $E_u$
exactly~\eqref{mm:remainder} and
\begin{align*}
 P_u&=\mathbf1_{\psi_u=1}\sqrt D\,\widehat W(1)
 \sum_{d\ \mathrm{squarefree}\atop(d,S)=1}
 \frac{\beta_U(d)\nu(d)\chi_d(u)}{Nd}\rho_{T(u,d)},\\
 \rho_T&=\frac{\operatorname*{res}_{s=1}\zeta_K(s)}{\zeta_K(2)}
                   \prod_{p\in T}(1+(Np)^{-1})^{-1}.
\end{align*}
This does not impose positivity on $P_u$. Since $\beta_U$ is supported
on $Nd\le U^2\le D$, every inner scale $Y$ is at least one.
Four integrations by parts give
$|\widehat W(\sigma+it)|\ll_\sigma\mathcal M_4(W)(1+|t|)^{-4}$,
which justifies the shifted integrals and their horizontal limits with
the classical Hecke growth estimates.

Apply Minkowski to~\eqref{mm:remainder} in $\ell^2$ of the actual
$u$ rows, then use~\eqref{mm:mask-bound} and the off-central moment.
For sufficiently small internal losses this bounds $\|E\|_2$ by
\[
 C D^\delta H^{1/2+O(\eta)}
 \sum_d\frac{|\beta_U(d)|}{(Nd)^{\sigma-\eta}}
 \int_{\mathbb R}|\widehat W(\sigma+it)|
                     (1+|t|)^{(1+O(\eta))/2}\,dt.
\]
All original $\nu(d)\chi_d(u)$ have modulus at most one and remain
inside the exact expression before this inequality. The actual truncated
factorization gives, without replacing $\beta_U$ by arbitrary coefficients,
\[
 \sum_d\frac{|\beta_U(d)|}{(Nd)^{\sigma-\eta}}
 \le\left(\sum_{Na\le U}(Na)^{-\sigma+\eta}\right)^2
 \ll_{\delta,\eta}U^{1-2\delta+2\eta}.
\]
Choose $\eta<\delta$ if necessary and use the fixed-field ideal count.
The height integral is bounded by $C\mathcal M_4(W)$. Squaring and
absorbing the small losses in $(HD)^\varepsilon$, using $U^2\le D$,
proves~\eqref{mm:energy}. Thus the formerly conditional numerator
moment input in this residual transfer is discharged, while its
positive-power cutoff loss is retained explicitly.

\subsection{What remains open and what has been verified}
The new remainder envelope improves the old scalar-convexity $H$ power
$3/2-\delta+\varepsilon$ to $1+\varepsilon$. At fixed $U$ one may
choose $\delta$ for a requested small-loss exponent; for $U=D^a$ with
fixed $a>0$, the remaining excess tends to $2a$ as $\delta\downarrow0$.
Neither statement controls $P_u$, the other signed piece, nor their
covariance. In the original Vaughan identity the short $2a_U$ term
must be restored at smaller scales where its support can meet the
annulus. Uniform $\mathcal M_4$ bounds do not justify interchanging a
profile supremum with the row moment or assuming $U^2\le D'$ after a
scale $D'$ is reduced.

The independent prose audit is recorded in
\path{planning/analytic-oracle/arithmetic-review/adversarial-afm-review-20261104.md}.
The source-hashed finite checker tests valuation/support decompositions,
squarefull coordinates, height exponents and rational Poisson envelopes.
Its finite success does not prove reciprocity, analytic continuation,
the large sieve, the approximate functional equation or the moments.
No new Lean declarations accompany these analytic theorems. Existing
Lean local Euler identities are a separately limited verification layer.
The independent positive arithmetic determinant comparison, original
signed low-gcd short-row estimate, all required contour uniformities and
GRH remain open; Dirichlet, Hecke, Dedekind and general Artin scopes are
not identified with one another.

\paragraph{Pinned primary sources and external inputs.}
The principal analytic family input is
\href{https://arxiv.org/pdf/1112.1650v1}{Blomer--Goldmakher--Louvel,
arXiv:1112.1650v1}, Theorem~1.3, Lemma~2.1, Appendix~II and Section~5.
Our fixed-$\nu$ proof changes column coefficients and redoes the
functional-equation reduction; it does not claim the published
corollary directly states every fixed twist.
For comparison,
\href{https://arxiv.org/pdf/2201.01885v2}{Gao--Zhao,
arXiv:2201.01885v2}, Lemma~2.9 and Section~3.2, uses the sextic
large-sieve/moment interface. Its GRH-conditional central-density theorem
is not invoked. Abelian Kummer identification uses the classical
reciprocity theorem, with the project's pinned Milne \emph{Class Field
Theory}, version~4.03 (6 August 2020), Theorem~0.8 and Chapter~I,
Section~1, a reference:
\href{https://www.jmilne.org/math/CourseNotes/CFT.pdf}{Milne's notes}.
The primitive degree-one Hecke continuation and completed functional
equation are explicit classical external inputs; the project's pinned
\href{https://u.cs.biu.ac.il/~reznikov/courses/kudla-1.pdf}{Kudla notes
on Tate's thesis}, Theorem~4.3 and Corollary~4.4, supply a reference.
Their sketch and incomplete displayed epsilon formulas are not used as
independently certified constants. Standard Stirling, Mellin inversion,
Poisson representation and scalar convexity complete the external
analytic dependencies. None is replaced by a finite diagnostic or a
GRH-equivalent assumption.


\chapter{Weighted cutoffs and original-family moments}
% Generated shared reviewed material.
\section{Recombining the actual weighted cutoff}
\label{rc:core}
Keep the finite-order Eisenstein family and original masks of
Section~\ref{mm:core}. For $U\ge1$ and a squarefree ideal $m$, put
\[
 A_U(m)=\sum_{d\mid m}\beta_U(d)
 =\sum_{ab\mid m\atop(a,b)=1,\ Na,Nb\le U}\mu(a)\mu(b).
\]
Every summation variable here is an integral ideal. The last formula
uses squarefreeness of $m$; an unconstrained square of a M\"obius
polynomial would also include forbidden overlapping supports.
\begin{theorem}[Exact recombination and an absolute-half-plane rate]
\label{rc:recombination}
Write $T_0=S\cup\operatorname{supp}(u)$ and
$J_{u,U}(s)=\sum_{d\ \mathrm{squarefree},\ (d,S)=1}
\beta_U(d)\nu(d)\chi_d(u)(Nd)^{-s}F_{u,d}(s)$.
For $\operatorname{Re}s>1$,
\begin{align*}
 J_{u,U}(s)&=\sum_{m\ \mathrm{squarefree},\ (m,T_0)=1}
                   A_U(m)\psi_u(m)(Nm)^{-s},\\
 |A_U(m)|&\le3^{\omega(m)},\qquad
 A_U(m)=\mu(m)\quad(Nm\le U).
\end{align*}
For every $\kappa>0$ and $0<\epsilon<\kappa$, uniformly for
$\operatorname{Re}s\ge1+\kappa$, all original rows and all $U\ge1$,
\[
 \left|J_{u,U}(s)-\frac1{L^{T_0}(s,\psi_u)}\right|
 \le U^{-\kappa+\epsilon}
 \{\zeta_K(1+\epsilon)^3+\zeta_K(1+\epsilon)\}.
\]
Here $L^{T_0}$ removes every prime of $T_0$, including conductor-drop
primes; the bound is in the absolute-convergence half-plane only.
\end{theorem}
\begin{proof}
On a surviving term $m=dc$, the two squarefree ideals $c,d$ are
coprime, and both avoid $T_0$. The original product
$\nu(d)\chi_d(u)\psi_u(c)$ is then $\psi_u(m)$. Grouping by $m$
gives $A_U$. Each prime of $m$ has three possible positions in
$(a,b,m/(ab))$, so the absolute sum is at most $3^{\omega(m)}$.
When $Nm\le U$ all divisors fit both cutoffs; the primewise sum is
$1-1-1=-1$, giving $\mu(m)$.

The same bound justifies absolute convergence and regrouping for
$\operatorname{Re}s>1$, because on squarefree ideals
$3^{\omega(m)}=\tau_3(m)$ and
$\sum_m\tau_3(m)(Nm)^{-v}=\zeta_K(v)^3$ for $v>1$,
where this latter sum runs over all nonzero integral ideals.
Subtract the absolutely convergent M\"obius series for $1/L^{T_0}$.
Only $Nm>U$ remain. Rankin's inequality with exponent
$\eta=\kappa-\epsilon$ bounds their absolute sum by
\[
 U^{-\eta}\sum_m\{\tau_3(m)+1\}
                  (Nm)^{-1-\epsilon},
\]
which proves the stated constant independently of deleted primes,
row multiplicities or character phases.
\end{proof}

For $W$ supported in $[1,2]$, the same exact regrouping gives the
original normalized type-I polynomial
\[
 I_u=D^{-1/2}\sum_{m\ \mathrm{squarefree},\ (m,S)=1}
 A_U(m)\nu(m)\chi_m(u)W(Nm/D).
\]
This formula retains all original zero sentinels. It permits a direct
family large sieve before any absolute $d$ sum; its quantitative moment
is developed separately below. The displayed absolute-half-plane rate
alone is not that critical-strip moment.

% Generated shared reviewed material.
\section{Euler packaging, coprimality and the analytic boundary}
\label{rc:details}
For $p\notin T_0$, set $z_p=\psi_u(p)(Np)^{-s}$ and
$t_p=z_p/(1+z_p)$. On $\operatorname{Re}s>1$, deleting an additional
surviving prime divides the squarefree Euler product by $1+z_p$.
Consequently the exact retained expression is
\[
 J_{u,U}(s)=F_{u,1}(s)B_{u,U}(s),\qquad
 B_{u,U}(s)=\sum_{d\ \mathrm{squarefree},\ (d,T_0)=1}
                 \beta_U(d)\prod_{p\mid d}t_p.
\]
Original zeros enforce $(d,u)=1$ even when a sixth power drops from
its primitive conductor. This factorization cannot restore a forbidden
$d$ row by replacing a zero sentinel with an unramified primitive value.
The equality continues wherever both sides are holomorphic.

Without the cutoffs, $\beta_\infty(d)=\mu(d)2^{\omega(d)}$ on
squarefree ideals. Absolute convergence justifies
\[
 B_{u,\infty}=\prod_{p\notin T_0}(1-2t_p)
 =\prod_{p\notin T_0}\frac{1-z_p}{1+z_p},\qquad
 F_{u,1}B_{u,\infty}=\prod_{p\notin T_0}(1-z_p)
                   =\frac1{L^{T_0}(s,\psi_u)}.
\]
Theorem~\ref{rc:recombination} gives a quantitative rate rather than an
unsupported exchange of this limit with a critical-strip contour.

\paragraph{A local counterexample to the naive square.}
In a one-prime projection at a good Eisenstein prime of norm $7$, take
the principal local phase and $s=1$. Then $z=1/7$, $t=1/8$ and,
for $U\ge7$, the exact polynomial is $B_U=1-2t=3/4$.
The unconstrained square $(1-t)^2=49/64$ differs by $1/64$:
it admits the forbidden pair $a=b=p$. After multiplying by
$F=1+z=8/7$ it gives $7/8$ instead of the correct $6/7$.
This is an actual local Euler calculation, not an infinite-series or
mean-square counterexample. It identifies the exact coprimality condition
that a proposed mollifier-square simplification must preserve.

\paragraph{The exact intermediate coefficient band.}
For a squarefree $m$ with $U<Nm\le U^2$, one has $A_U(m)=-\mu(m)$.
Indeed the unrestricted disjoint-assignment sum is $\mu(m)$. If $Na>U$,
every $b\mid m/a$ automatically has $Nb\le Nm/Na<U$, and
$\sum_{b\mid m/a}\mu(b)$ vanishes unless $a=m$. Hence the excluded
$a$ assignments contribute exactly $\mu(m)$; the excluded $b$
assignments contribute another $\mu(m)$. They cannot overlap because
$NaNb>U^2\ge Nm$ contradicts $ab\mid m$. Subtraction proves the claim.
Together with the small-norm band, this gives $A_U=\mu$ for $Nm\le U$
and $A_U=-\mu$ for $U<Nm\le U^2$. The original smoothed columns have
$Nm\ge D\ge U^2$, so the new band does not justify replacing their
coefficients by $-\mu(m)$: generally they lie beyond it.

\begin{theorem}[A stronger cutoff convergence would require zero-freeness]
\label{rc:boundary}
Fix a nonprincipal primitive $\psi_u$. If these actual $J_{u,U}$
converge locally uniformly on $\operatorname{Re}s>1/2$ as
$U\longrightarrow\infty$, then $L(s,\psi_u)$ has no zero there.
\end{theorem}
\begin{proof}
Classical abelian Hecke continuation makes the numerator of
$F_{u,1}=L^{T_0}(s,\psi_u)/L^{T_0}(2s,\psi_{u,2})$ holomorphic.
Its denominator is nonzero on $\operatorname{Re}s>1/2$ by the
absolute Euler product at $2s$. Each $B_{u,U}$ is a finite sum with
nonzero denominators $1+z_p$ on $\operatorname{Re}s>0$.
Thus every $J_{u,U}$ is holomorphic on the asserted half-plane.
A locally uniform limit $G$ is holomorphic there. On
$\operatorname{Re}s>1$, the already proved convergence gives
$L^{T_0}(s,\psi_u)G(s)=1$. The identity theorem extends this equality
to the connected half-plane, precluding a zero of $L^{T_0}$.
The finitely deleted factors $1-\psi_u(p)(Np)^{-s}$ do not vanish
when $\operatorname{Re}s>0$, so the same conclusion holds for the
primitive $L$-function.
\end{proof}
This is an implication, not a proof that the stronger convergence holds,
and it does not assert a converse. In particular, a family-averaged energy
bound is not silently promoted to locally uniform convergence for each
individual character. Principal rows have a pole at $1$ and were excluded
from this proposition. General Artin holomorphy is not an input.

\paragraph{Reproducible exact diagnostics.}
The new source \path{examples/arithmetic_cutoff/exact_recombination.py}
uses actual split/inert prime ideals and raw sextic symbols in
$\mathbb Z[\omega]$, with exact rational arithmetic. It checks the direct
masked sum, coefficient regrouping, Euler packaging, both finite full-cutoff
limits and original zeros, including third/sixth-power rows. Its overlapping-
support negative controls are not tests of an infinite large sieve or of
zero location. The analytic convergence argument above uses ordinary ideal
Euler products and the explicitly stated Hecke continuation theorem; no
new Lean declaration or complete GRH proof is claimed.

% Generated shared reviewed material.
% Original masked cutoff family; no new Lean or GRH theorem.
\section{A cutoff-uniform envelope for the original masked remainder}
\label{wc:core}
Fix $K=\mathbb Q(\omega)$, $\mathcal O=\mathbb Z[\omega]$, a finite-order
Hecke character $\nu$, and a finite set $S$ containing the primes above $6$
and the conductor of $\nu$. All unspecified ideal sums below are over
nonzero integral ideals. Let $\chi_m(u)=(u/m)_6$ be the original symbol,
zero when $(m,u)>1$, and put
\[
 \beta_U(d)=\sum_{ab=d\atop Na,Nb\le U}\mu(a)\mu(b),\qquad
 A_U(m)=\sum_{d\mid m}\beta_U(d).
\]
The columns $d,c,m$ are squarefree and prime to $S$ where indicated.
Retain the primitive characters $\psi_u=\operatorname{prim}(\nu\kappa_u)$,
$\psi_{u,2}=\operatorname{prim}(\psi_u^2)$, and the full deleted set
$T(u,d)=S\cup\operatorname{supp}(u)\cup\operatorname{supp}(d)$.
For $H,D,U\ge1$, $U^2\le D$, and $W\in C_c^\infty((0,\infty))$
supported in $[1,2]$, the original normalized type-I piece is
\begin{equation}\label{wc:polynomial}
 I_u=D^{-1/2}\sum_{m\ \mathrm{squarefree}\atop(m,S)=1}
     A_U(m)\nu(m)\chi_m(u)W(Nm/D).
\end{equation}
Its exact principal residue is
\begin{align}\label{wc:principal}
 P_u&=\mathbf1_{\psi_u=1}\sqrt D\,\widehat W(1)
 \sum_{d\ \mathrm{squarefree}\atop(d,S)=1}
 \frac{\beta_U(d)\nu(d)\chi_d(u)}{Nd}\rho_{T(u,d)},\notag\\
 \rho_T&=\frac{\operatorname*{res}_{s=1}\zeta_K(s)}{\zeta_K(2)}
               \prod_{p\in T}(1+(Np)^{-1})^{-1}.
\end{align}
Define $E_u=I_u-P_u$; this is the same residue-subtracted remainder as the
fully masked Mellin contour, with no cutoff limit or profile supremum.

\begin{theorem}[Uniformity in the actual truncated cutoff]
\label{wc:uniform}
For every $\varepsilon>0$,
\begin{equation}\label{wc:bound}
 \sum_{0<Nu\le H}|E_u|^2
 \ll_{\nu,S,\varepsilon}\|W\|_\infty^2(HD)^\varepsilon
       \{H+DH^{1/3}+(HD)^{2/3}\}.
\end{equation}
The sum is over every actual nonzero Eisenstein element, including units,
$S$-supported powers and sixth-power rows. The constant is independent of
$U$ under $U^2\le D$. This is an application of the classical sextic
large sieve with exact original coefficients, not a zero-location theorem.
\end{theorem}
\begin{proof}
On an absolutely convergent line, expand the restored squarefree series
$F_{u,d}$. Its condition $(c,T(u,d))=1$ includes $(c,d)=1$, and the
original symbols multiply. Grouping $m=cd$ gives~\eqref{wc:polynomial}
by Mellin inversion. Each prime of $m$ may belong to the first truncated
M\"obius divisor, to the second, or to neither. Consequently
$|A_U(m)|\le3^{\omega(m)}\ll_\eta(Nm)^\eta$, uniformly in $U$.
The length-$D$ coefficient sequence in~\eqref{wc:polynomial} therefore
has squared norm $\ll_\eta\|W\|_\infty^2D^\eta$.

The all-ideal row norm in Section~3 of
\href{https://arxiv.org/pdf/1112.1650v1}{Blomer--Goldmakher--Louvel,
arXiv:1112.1650v1} gives
\[
 B_2^1(X,D)\ll_\eta(XD)^\eta
                 \{X+DX^{1/3}+(XD)^{2/3}\}.
\]
For completeness, split $a=a_1a_2^2\cdots a_5^5a_6^6$ dyadically;
$a_1,\ldots,a_5$ are squarefree and pairwise coprime. If their dyadic
sizes are $A_j$, put $R=\prod_{j=1}^6A_j$ and $M=\max(A_1,A_2)$.
Pivot on the first or second variable, using the squarefree powered
sieve proved simultaneously in that paper's Section~3.2. The other
variables are bounded column multipliers. The block energy is at most
\[
 (XD)^\eta\{R+DR/M+D^{2/3}R/M^{1/3}\}\sum_m|\lambda_m|^2.
\]
The inequalities $R\le X$, $R^3\le XM^3$ and $R^3\le X^2M$
follow from $\prod_jA_j^j\le X$ and the pivot choice. Summing the
logarithmically many blocks proves the displayed bound.

Finite $S$-unit and ray-class partitions convert the original literal
symbol into the canonical column value $\vartheta_m(a)$ with fixed
modulus-one coefficient multipliers. On noncoprime pairs both values
are zero: squarefree $m$ has outside-$S$ conductor exactly $m$.
This orientation preserves sixth-power zeros. At most six elements
generate each outside-$S$ ideal times an $S$-supported ideal; the
$S$-supported sums converge at powers $1$, $1/3$ and $2/3$.
Thus~\eqref{wc:bound} holds first with $I_u$ in place of $E_u$.

A principal $\psi_u$ requires every outside-$S$ valuation of $u$ to
be divisible by $6$. There are $O_{\nu,S}(H^{1/6})$ such literal
rows. Equation~\eqref{wc:principal}, $|\rho_T|\ll_K1$, and the actual
coefficient sum
\[
 \sum_d\frac{|\beta_U(d)|}{Nd}
 \le\left(\sum_{Na\le U}\frac1{Na}\right)^2\ll\log^2(2U)
\]
give $\sum|P_u|^2\ll DH^{1/6}\|W\|_\infty^2\log^4(2U)$.
This is absorbed by the right side of~\eqref{wc:bound}, using
$U^2\le D$ and a smaller internal exponent. Finally
$|E_u|^2\le2|I_u|^2+2|P_u|^2$. The extended proof specifies the
canonical-family transfer and the powered-sieve input precisely.
\end{proof}

Combining this with the previously proved fixed-twist AFM transfer yields
\begin{equation}\label{wc:minimum}
 \sum_{0<Nu\le H}|E_u|^2\ll_{\nu,S,\delta,\varepsilon}
 \mathcal M_4(W)^2(HD)^\varepsilon
 \min\{D^{2\delta}U^{2-4\delta}H,\ H+DH^{1/3}+(HD)^{2/3}\}.
\end{equation}
At $H=D^{3/4}$ and $U=D^a$, $0\le a\le1/2$, the excess over $H$
is at most $\min\{2\delta+(2-4\delta)a,1/2\}$, up to arbitrarily
small losses. The new envelope improves the old one for $a>1/4$ as
$\delta\downarrow0$, but still leaves an exponent loss of $1/2$.
It does not settle the original signed short-row bound, type-II covariance,
positive arithmetic determinant comparison or GRH. The principal-residue
bound used here is sufficient for this envelope, not the stronger target
needed to reassemble the full signed argument.

% Generated shared reviewed material.
\section{All-ideal reciprocity and the surviving block obstruction}
\label{wc:details}
The proof uses no new zero-free input. Its external analytic ingredient
is the pinned BGL large sieve and the powered estimates established in
its proof; elementary ideal counting and the already restored Mellin
identity supply the remaining interfaces. All constants may depend on
fixed $K$, $S$, $\nu$ and declared small exponents, never on $U$.

\subsection{Exact recombination retains the primitive corrections}
On $\operatorname{Re}s>1$, start from the full defining series, with
its original zero masks:
\[
 F_{u,d}(s)=\sum_{c\ \mathrm{squarefree}\atop(c,T(u,d))=1}
                 \nu(c)\chi_c(u)(Nc)^{-s}.
\]
Its primitive ratio involves $L(s,\psi_u)$, the independently primitive
$L(2s,\psi_{u,2})$ and every deleted Euler factor. Expanding this
\emph{restored} series gives the exact equality
\begin{align}\label{wc:recombine}
 &\sum_{d\ \mathrm{squarefree}\atop(d,S)=1}
 \beta_U(d)\nu(d)\chi_d(u)(Nd)^{-s}F_{u,d}(s)\notag\\
 &\qquad=\sum_{m\ \mathrm{squarefree}\atop(m,S)=1}
               A_U(m)\nu(m)\chi_m(u)(Nm)^{-s}.
\end{align}
Indeed $(c,d)=1$ makes $m=cd$ squarefree, with $d\mid m$. If $m$
meets $u$, both original symbols retain their zero. Only on the
absolutely convergent line is this infinite series used. Mellin inversion
there gives the finite smoothed polynomial~\eqref{wc:polynomial}.
Nothing evaluates~\eqref{wc:recombine} as a Dirichlet series on
$\operatorname{Re}s=1/2+\delta$ or interchanges a cutoff limit with
continuation across unknown zeros.

For squarefree $m$, a factorization $ab=d\mid m$ assigns each prime
to $a$, to $b$, or to neither. The actual inequalities $Na,Nb\le U$
select a subset of these assignments. Thus
\[
 |A_U(m)|\le3^{\omega(m)}\ll_\eta(Nm)^\eta.
\]
The last divisor bound follows by choosing a fixed threshold so
$3\le(Np)^\eta$ beyond it, then absorbing the finitely many smaller
prime ideals in a constant. When $Nm\le U$ all assignments occur,
so $A_U(m)=\prod_{p\mid m}(1-1-1)=\mu(m)$; no such simplification
is made for the length-$D$ columns in general. With
$\lambda_m=D^{-1/2}A_U(m)\nu(m)W(Nm/D)$, ideal counting and a
smaller internal $\eta$ give
\begin{equation}\label{wc:l2}
 \sum_m|\lambda_m|^2\ll_\eta\|W\|_\infty^2D^\eta,
 \qquad \operatorname{supp}\lambda\subseteq\{D\le Nm\le2D\}.
\end{equation}

\subsection{The powered squarefree input and all-ideal pivot}
Let $\vartheta_a$ be BGL's canonical primitive sixth-order family.
Their Section~3 defines $B_1^j$ using squarefree row/column ideals and
kernel $\vartheta_a(m)^j$, while $B_2^j$ uses \emph{all} ideal rows
and kernel $\vartheta_m(a)^j$. Theorem~1.3 states the squarefree bound
for $j=1$; its Section~3.2 proves the estimates simultaneously for
$1\le j\le5$, because the recursion in Lemma~3.4 uses $j\ell$.
In particular the two inputs needed here are
\begin{equation}\label{wc:powered-input}
 B_1^j(M,D)\ll_\eta(MD)^\eta
                 \{M+D+(MD)^{2/3}\},\qquad j=1,2.
\end{equation}
To see the final step in the published recursion explicitly, its
$E_{4/3}$ estimate is $M^{4/3}+D+(MD)^{2/3}$.
If $M\le D$, its first term is at most $(MD)^{2/3}$; if $M>D$,
use reciprocity/duality (3.2) and the estimate with $M,D$ reversed.
The finite set of powers is treated together; no inference from a
fixed conductor count to an arbitrary primitive-character sieve occurs.

Write each ideal row uniquely as
\[
 a=a_1a_2^2a_3^3a_4^4a_5^5a_6^6,
\]
where $a_1,\ldots,a_5$ are pairwise-coprime squarefree ideals and
$a_6$ is arbitrary. For valuation $6q+r$, $0\le r\le5$, place the
prime in $a_r$ if $r>0$ and its power $q$ in $a_6$.
Dyadically let $A_j\le Na_j<2A_j$, with $A_j\ge1$, and put
$R=\prod_{j=1}^6A_j$. A nonempty block satisfies
$\prod_jA_j^j\le X$ when $Na\le X$.
Choose $\ell=1$ if $A_1\ge A_2$, and $\ell=2$ otherwise; put
$M=A_\ell$. Keep every other ideal fixed. Evaluation of the single
\emph{column} character gives
\[
 \vartheta_m(a)=\vartheta_m(a_\ell)^\ell
                  \vartheta_m(a/a_\ell^\ell).
\]
The second factor is a bounded column multiplier, including its original
zero if $m$ meets any fixed factor. It may depend on all the fixed
ideals but not on the pivot row. Pairwise-coprimality and the exact row
norm restriction merely select a subset of the positive pivot-row sum.
Apply~\eqref{wc:powered-input} to the first factor. There are
$O_K(R/M)$ choices of the other ideals by ideal counting. The block
energy is bounded by
\begin{equation}\label{wc:block}
 (XD)^\eta\{R+DR/M+D^{2/3}R/M^{1/3}\}\sum_m|\lambda_m|^2.
\end{equation}

All three monomial bounds used to sum these blocks are elementary.
First $R\le\prod_jA_j^j\le X$.
If $M=A_1$, then
\[
 (R/M)^3=A_2^3\cdots A_6^3
 \le A_1A_2^2A_3^3A_4^4A_5^5A_6^6\le X,
\]
since $A_2\le A_1$ and $j\ge3$ for the remaining factors.
If $M=A_2$, use $A_1\le A_2$ in the same comparison.
Finally,
\[
 \frac{R^3}{M(\prod_jA_j^j)^2}
 =\frac{A_1}{MA_2A_3^3A_4^5A_5^7A_6^9}\le1.
\]
Thus $R/M\le X^{1/3}$ and $R/M^{1/3}\le X^{2/3}$.
Insert these in~\eqref{wc:block} and sum the logarithmically many
blocks. This proves the all-ideal estimate used in Theorem~\ref{wc:uniform}.
It is precisely the mechanism of BGL Lemma~3.4 and (3.6)--(3.9),
with the final powered bound inserted. The simpler fixed-squarefull-row
argument would leave the weaker $D\sqrt X$ term.

\subsection{Literal rows, reciprocal orientation and principal count}
Let $a=(u)_S$ denote the outside-$S$ ideal. The finite $S$-unit
comparison gives $\kappa_u=\eta_u\vartheta_a$ as primitive
characters, with $\eta_u$ in a finite fixed set.
For $(a,m)=1$, BGL reciprocity therefore gives
\begin{equation}\label{wc:literal-reciprocity}
 \chi_m(u)=\eta_u(m)C([a],[m])\vartheta_m(a).
\end{equation}
The notation is deliberate: the left side is the project's original
$(u/m)_6$, not BGL's canonical character indexed by $m$.
When $(a,m)>1$, the left side is zero. The right side is also zero,
because squarefree $m$ has outside-$S$ primitive conductor exactly $m$.
Equation~\eqref{wc:literal-reciprocity} hence extends to every pair.
Using $\vartheta_a(m)$ on noncoprime pairs instead would be false:
primitive $\vartheta_{b^6}$ has restored primes of $b$, whereas the
original symbol must still vanish there.

Partition rows by $\eta_u$ and $[a]$, both finite. On each part the
factor $\eta_u(m)C([a],[m])$ is a fixed modulus-one column multiplier.
A positive row sum over that part is bounded by the full $B_2^1$ row
sum. For each $S$-supported ideal $b$ and outside-$S$ ideal $a$, at
most six actual elements generate $ab$. Thus the required literal energy
is bounded by the all-ideal norm summed over $b$, with $X=H/Nb$.
The geometric products
\[
 \sum_{b\ S\text{-supported}}(Nb)^{-q}
   =\prod_{p\in S}(1-(Np)^{-q})^{-1}<\infty,
 \qquad q=1,\ 1/3,\ 2/3,
\]
transfer the three terms without losing any unit or $S$-power row.
Together with~\eqref{wc:l2}, this proves the claimed $I$ energy.

If $\psi_u$ is principal, $\nu$ is unramified outside $S$, and tame
Kummer inertia at $v_p(u)$ has order $6/\gcd(6,v_p(u))$.
It must be trivial, so each such valuation is divisible by $6$.
Writing $(u)=ba_6^6$ with $b$ supported in $S$, ideal counting gives
\[
 \#\{0<Nu\le H:\psi_u=1\}
 \ll_K H^{1/6}\prod_{p\in S}(1-(Np)^{-1/6})^{-1}.
\]
This is only a necessary condition, sufficient for an upper count.
The exact residue formula~\eqref{wc:principal} and
$|\widehat W(1)|\le\|W\|_\infty$ then give the principal energy
bound used above. No sign or positivity of that residue is assumed.

\subsection{What the envelope still fails to prove}
The $DH^{1/3}$ term is still larger than $H(HD)^\varepsilon$ at
$H=D^{3/4}$. The monomial comparison itself is saturated by
$A_1=A_2=X^{1/3}$ and $A_3=\cdots=A_6=1$.
This proves sharpness of that \emph{particular comparison}, not an
energy lower bound for actual arithmetic coefficients. In particular,
no conclusion that the signed $A_U$ cannot cancel is justified.
The next load-bearing input must exploit those actual truncated M\"obius
coefficients inside the mixed moment, or establish the required signed
type-II minus type-I covariance. The smaller-scale short term, principal
contribution, full arithmetic determinant identification and GRH remain
separate obligations.

Exact finite checks use the split Eisenstein primes generated by
$3+\omega$, $4+\omega$, $5+2\omega$, $6+\omega$, of norms
$7,13,19,31$. Their actual residue symbols are computed in
$\mathbb Q(\zeta_6)$, retaining sixth-power zeros. Coefficient
recombination and dyadic inequalities pass; they check finite algebra
and bookkeeping, not the external analytic sieve or the contour theorem.
No new Lean declarations accompany this argument.


% Generated shared reviewed material.
\section{High cutoffs, actual bilinear factors and a signed covariance test}\label{sb:core}
Retain the original Eisenstein family, fixed finite-order $\nu$, fixed $S$,
common $W\in C_c^\infty((0,\infty))$ supported in $[1,2]$, and every raw
shared-prime zero. Put
\[
 C_U(m)=\sum_{\substack{abc=m\\Na>U,\,Nb>U}}\mu(a)\mu(b),\qquad
 r_U(m)=\#\{(a,b):ab=m,\ Na,Nb>U\}.
\]
All factors are ideals, squarefree and pairwise coprime; on the active
annulus $A_U=C_U-\mu$. A coefficient identity is different from a row-energy bound.

\begin{proposition}[The high-cutoff window]\label{sb:window}
If $D>1$ and $2D/7\le U^2\le D$, then for every surviving squarefree
column $D\le Nm\le2D$,
\[
 C_U(m)=\mu(m)r_U(m),\qquad
 A_U(m)=\mu(m)(r_U(m)-1),\qquad |A_U(m)|\ge1.       
\]
For each fixed nonzero $W$, the normalized squared coefficient norm satisfies
\begin{equation}
 D^{-1}\sum_{\substack{m\ {\rm sf}\\(m,S)=1}}
 |A_U(m)W(Nm/D)|^2\gg_{W,S}1                     \label{sb:coefficient-lower}
\end{equation}
for large $D$, uniformly in this window. This is not a lower bound for row energy.
\end{proposition}
\begin{proof}
The least norm of a nonunit ideal outside $S$ is at least $7$. A long-long
term with $c\ne1$ would have $Nm>7U^2\ge2D$. Thus $ab=m$, including both
endpoints. The fixed-point-free swap $(a,b)\mapsto(b,a)$ makes $r_U$ even,
so $A_U$ is an odd nonzero integer. The last claim follows by counting
squarefree ideals on a fixed interval where $|W|$ is bounded below; the
constant and error are proved in the extended proof.
\end{proof}

The exact remaining polynomial is genuinely bilinear, with $(a,b)=1$,
$U<Na,Nb<2D/U$, $D\le Na\,Nb\le2D$, and
$x_{ab}(u)=x_a(u)x_b(u)$, including zeros. Both factors see the same row
$u=\varepsilon\alpha\beta^2$. Cauchy followed by the accepted powered BGL
bound gives only
\begin{equation}
 Q_{B_X}(C_U,C_U)\ll_{\nu,S,\epsilon}
 (DX)^\epsilon\|W\|_\infty^2
 \{\sqrt D\,X^2+DX+D^{5/6}X^{5/3}\}.              \label{sb:factor-bound}
\end{equation}
At $X\asymp D^{1/4}$ the exponents are $1,5/4,5/4$; the earlier regrouped
bound has $1/2,5/4,13/12$ and is stronger. No $BT_\sigma$ saving is obtained.
Multiplying two marginal operator norms does not supply a coupled fourth moment.

\paragraph{A finite test in the original family.}
There is no automatic nonnegative signed covariance, even when
$C_U(m)\mu(m)\ge0$. Take $\nu=1$, $S$ the primes above $2,3$, $D=400$,
$U=17$, and all $72$ literal rows of $B_7$. Choose the common nonnegative
smooth profile constructed in the extended proof, with values $10$ at $439/400$, $1$ at
$475/400$, and zero at every other integer norm divided by $400$.
There are exactly four active squarefree ideal columns, of norms
$439,439,475,475$. In that order,
\[
 A=(1,1,1,1),\quad C=(0,0,2,2),\quad
 \mu=(-1,-1,1,1),\qquad
 \Re G_{B_7}=
 \begin{pmatrix}72&0&0&18\\0&72&18&0\\0&18&72&0\\18&0&0&72\end{pmatrix}.
\]
Exact raw-symbol enumeration, independently reproducible from the ideal
and residue-field definitions, gives
\begin{equation}
 Q(C,\mu)=-\frac{27}{25},\qquad
 Q(C,C)=\frac{36}{25},\quad
 Q(\mu,\mu)=\frac{864}{25},\quad
 Q(A,A)=\frac{954}{25}.                          \label{sb:covariance-example}
\end{equation}
The real Gram matrix is positive definite, with eigenvalues $54,90$;
its positivity does not force a cross covariance to be positive.
This fixed finite example refutes that shortcut, not the asymptotic
$BT_\sigma$ target or the original signed mean-square estimate.

\paragraph{The two-cutoff contour remains coupled.}
For $T_0=S\cup\operatorname{supp}(u)$, $\psi=\psi_u$ and independently
primitive $\psi_2=\operatorname{prim}(\psi^2)$, put $q_p=\psi(p)(Np)^{-s}$,
$\alpha_p=(Np)^{-z}$, $\beta_p=(Np)^{-w}$ outside $T_0$. The actual
three-placement Euler factor is $1+q_p-q_p\alpha_p-q_p\beta_p$. Its exact
factorization is
\begin{equation}
 \mathcal F(s,z,w)=
 \frac{L^{T_0}(s,\psi)}{L^{T_0}(s+z,\psi)L^{T_0}(s+w,\psi)}
 \prod_{p\notin T_0}\left(1-
 \frac{q_p^2(1-\alpha_p)(1-\beta_p)}
 {(1-q_p\alpha_p)(1-q_p\beta_p)}\right).           \label{sb:perron-factor}
\end{equation}
Exact double Perron inversion is valid on positive auxiliary lines
$c,d>1-\Re s$ for $1/2<\Re s<1$, retaining the original principal pole
exclusion. These safe lines reproduce the previous $U^{2-4\delta}$
energy loss at $\Re s=1/2+\delta$, with the same common-profile derivative
and fixed-margin dependencies. The correction product alone does not
justify moving those lines through reciprocal-$L$ poles. In fact
$\mathcal F(s,z,0)=1/L^{T_0}(s+z,\psi)$ as a meromorphic identity.
No zero-free region, reciprocal moment, analytic saving, new Lean
formalization or GRH proof is supplied by this construction.

% Generated shared reviewed material.
\section{Proof details for the high-window and two-cutoff tests}\label{sb:details}
\subsection{The coefficient lower bound and its precise limitation}
The element lattice $\mathbb Z[\omega]\subset\mathbb C$ has covolume
$\sqrt3/2$. Counting disks with boundary error $O(\sqrt x)$ and dividing
by the six units, using that this ring is a PID, gives
\[
 N_K(x)=\kappa x+O(\sqrt x),\qquad \kappa=\frac{\pi}{3\sqrt3}.
\]
Finite inclusion-exclusion at $S$ gives
$N_{K,S}(y)=\kappa y\prod_{p\in S}(1-(Np)^{-1})+O_S(\sqrt y)$.
Insert $1_{\rm sf}(m)=\sum_{d^2\mid m}\mu(d)$. Absolute convergence at
$2$, the tail $\sum_{Nd>\sqrt x}(Nd)^{-2}=O(x^{-1/2})$, and
$\sum_{Nd\le\sqrt x}(Nd)^{-1}=O(\log(2x))$ yield
\[
 N_{{\rm sf},S}(x)=c_Sx+O_S(\sqrt x\log(2x)),\qquad
 c_S=\frac{\kappa}{\zeta_K(2)}\prod_{p\in S}(1+(Np)^{-1})^{-1}>0.
\]
These tail estimates follow by partial summation from the displayed lattice
count. For a fixed nonzero $W$, choose $[a,b]\subset(1,2)$ with
$|W|\ge\eta>0$. Proposition~\ref{sb:window} gives, for large $D$,
\[
 D^{-1}\sum|A_U(m)W(Nm/D)|^2\ge\tfrac12\eta^2c_S(b-a).
\]
There is no uniform lower constant over all profiles with one specified
supremum norm. Nor can a coefficient lower bound be combined with an
operator upper bound to produce a row-energy lower bound: no lower singular
value has been established on the actual vector.

\subsection{What legal factor Cauchy actually proves}
For dyadic $Na\asymp A,Nb\asymp B$, retain the exact $(a,b)=1$ mask,
hyperbolic support and $W(Na\,Nb/D)$ in
\[
 T_a(u)=\sum_b\mu(b)\nu(b)W(Na\,Nb/D)x_b(u).
\]
Cauchy gives
\[
 |\sum_a\mu(a)\nu(a)x_a(u)T_a(u)|^2\ll A\sum_a|T_a(u)|^2.
\]
For fixed $a$, fixed second row factor and fixed unit/ray classes, the
$b$ coefficients are independent of the pivot row and have squared norm
$\ll B\|W\|_\infty^2$. The prior classical power-one/two BGL input has
$\Delta(X,B)=X+B+(XB)^{2/3}$. Summing the $O(A)$ fixed columns and
$O(X)$ other row factors gives
\[
 \mathcal E_{A,B}\ll_{\nu,S,\epsilon}(DX)^\epsilon\|W\|_\infty^2
 \frac{A^2B}{D}X\Delta(X,B)
 \ll (DX)^\epsilon\|W\|_\infty^2
 \{AX^2+DX+D^{2/3}A^{1/3}X^{5/3}\}.
\]
In the high window $A,B\asymp\sqrt D$, and only boundedly many dyadic
pairs occur. This proves~\eqref{sb:factor-bound}; exchanging factors
or row pivots leaves its $DX$ term. A product of marginal second-moment
bounds supplies no estimate for the two transforms on the same $(r,s)$.
Any useful saving must control that coupled fourth moment or the signed
actual-coefficient covariance, with all masks retained.

\subsection{A smooth profile and the exact covariance computation}
Let
\[
 \phi(t)=\begin{cases}\exp(1-1/(1-t^2)),&|t|<1,\\0,&|t|\ge1,\end{cases}
 \qquad W_x(y)=x\phi(4(400y-439))+\phi(4(400y-475)).
\]
For $x\ge0$ this is nonnegative and smooth; its disjoint supports lie
strictly inside $(1,2)$. For every integer $n$, $W_x(n/400)$ equals
$x$ at $n=439$, $1$ at $n=475$ and zero otherwise. Thus the computation
uses exact values of a genuine common $C_c^\infty$ profile, rather than
claiming a finite rational weight is smooth. At $x=10$, $\|W_x\|_\infty=10$.

The primary ideal generators, in the order used above, are
\[
 (-23,-18),\quad(-5,18),\quad(10,-15),\quad(25,15),
\]
where $(a,b)$ denotes $a+b\omega$. The first two have norm $439$.
The last two factor as $(-2,3)(-5,0)$ and $(-5,-3)(-5,0)$, with norms
$19\cdot25$. Since $19,25>17$, the long-factor count is $2$ there and
zero on the prime columns. All squarefree ideals at both selected norms
are included. The complete $B_7$ uses the four good primes of norms
$7,7,13,13$ for distinct ordered row factors and all six units. It has
$72$ distinct rows, maximum norm $2197$; each is nonprincipal by its
outside-$S$ Kummer valuation $1$ or $2$ modulo $6$, whose nontrivial
tame character has conductor exponent $1$.

For the raw Gram $G_{ij}=\sum_{u\in B_7}x_{m_i}(u)
\overline{x_{m_j}(u)}$, exact finite-field exponentiation and lattice
enumeration give the displayed real matrix. The weighted real vectors are
$\ell=(0,0,2,2)$, $v=(-x,-x,1,1)$, and $a=(x,x,1,1)$.
Therefore
\[
 Q(C_U,\mu)=D^{-1}\ell^t(\Re G)v=\frac{36-9x}{50}.
\]
Its diagonal contribution is $36/50>0$ and its off-diagonal contribution
is $-9x/50$. For $x=10$ this is $-27/25$ and gives
\eqref{sb:covariance-example}. The identity
$(\Re G-72I)^2=18^2I$ proves the real eigenvalues $54,90$;
these are not being asserted as the eigenvalues of the full complex Gram.
The exact source is \path{examples/arithmetic_cutoff/signed_covariance.py}.
A second test retains raw zeros on the $246$ rows $0<Nu\le89$ whose
outside-$S$ part is nonunit. Its covariance at $x=10$ is $-69/50$.
This subset is guaranteed nonprincipal; it is not called the entire
nonprincipal family. Here $89=\lfloor400^{3/4}\rfloor$.
These are finite counterexamples to automatic covariance positivity,
not asymptotic counterexamples to $BT_\sigma$ or GRH.

\subsection{Double Perron without a divergent critical-strip series}
For $T_0=S\cup\operatorname{supp}(u)$ retain
$F_0(s)=L^{T_0}(s,\psi)/L^{T_0}(2s,\psi_2)$, with the independently
primitive square character. Set $t_p=q_p/(1+q_p)$ outside $T_0$.
The finite actual cutoff is
\[
 J_U(s,u)=F_0(s)B_U(s),\qquad
 B_U(s)=\sum_{\substack{a,b\ {\rm sf},\ (a,b)=1\\
                  (ab,T_0)=1,\ Na,Nb\le U}}
 \mu(a)\mu(b)\prod_{p\mid ab}t_p.
\]
For $\sigma=\Re s>1/2$, the series
\[
 \mathcal G(s,z,w)=\sum_{\substack{a,b\ {\rm sf},\ (a,b)=1\\(ab,T_0)=1}}
 \mu(a)\mu(b)\prod_{p\mid ab}t_p\,(Na)^{-z}(Nb)^{-w}
 =\prod_{p\notin T_0}(1-t_p(Np)^{-z}-t_p(Np)^{-w})
\]
converges absolutely on $\sigma+\Re z>1$, $\sigma+\Re w>1$, uniformly
in all vertical parts at fixed positive margins. This follows from
$|t_p|\le(Np)^{-\sigma}/(1-(Np)^{-\sigma})$. There
$\mathcal F=F_0\mathcal G$, without evaluating an infinite $c$ series
in the critical strip. Initial absolute convergence also proves
\eqref{sb:perron-factor} by the local identity
\[
 (1-q)(1+q-q\alpha-q\beta)
 =(1-q\alpha)(1-q\beta)-q^2(1-\alpha)(1-\beta).
\]
The correction product is holomorphic on
\[
 \sigma+\Re z>0,\quad\sigma+\Re w>0,\quad
 2\sigma+\min(\Re z,0)+\min(\Re w,0)>1.
\]
On margins $\sigma+\Re z,\sigma+\Re w\ge\eta>0$ and last exponent
$\ge1+\epsilon$, its local deviation from $1$ is at most
$4(1-2^{-\eta})^{-2}(Np)^{-1-\epsilon}$. This proves normal convergence
and a phase/deletion-uniform vertical upper bound. Nonvanishing of the
correction is not asserted.

For $1/2<\sigma<1$, positive lines $c,d>1-\sigma$, and
$V=\lfloor U\rfloor+1/2$, define
\[
 \mathcal I_{T_1,T_2}=\frac1{(2\pi i)^2}
 \int_{c-iT_1}^{c+iT_1}\int_{d-iT_2}^{d+iT_2}
 \mathcal F(s,z,w)\frac{V^{z+w}}{zw}\,dw\,dz.
\]
The half-integer $V$ includes the actual equality cutoff, since
$Na\le U$ is exactly $Na<V$. Put
$\Delta_V=\inf_{n\in\mathbb Z_{\ge1}}|\log(V/n)|\ge1/(2V+1)$, and let
$\mathcal A=\sum_{a,b}|\mu(a)\mu(b)\prod_{p\mid ab}t_p|
(Na)^{-c}(Nb)^{-d}$ with the same disjoint support.
For $c,T>0$ and $v>0$, $v\ne1$, define
\[
 K_{c,T}(v)=\frac1{2\pi i}\int_{c-iT}^{c+iT}\frac{v^z}{z}\,dz.
\]
Finite-height termwise integration of the absolutely convergent
$\mathcal G$ and the one-variable tail estimate
$|K_{c,T}(v)-1_{v>1}|\le2v^c/(\pi T|\log v|)$ give
\begin{equation}
 |\mathcal I_{T_1,T_2}-J_U|
 \le |F_0|\mathcal A V^{c+d}
 \left\{\frac2{\pi T_1\Delta_V}+\frac2{\pi T_2\Delta_V}
       +\frac4{\pi^2T_1T_2\Delta_V^2}\right\}.      \label{sb:perron-error}
\end{equation}
Thus $\mathcal I_{T_1,T_2}\to J_U$ as both heights tend independently to
infinity, away from the original principal pole. This is an oscillatory
limit, not an absolutely integrable infinite vertical integral.

Taking $T_1=T_2=V^2$ gives the safe upper bound
$|J_U|\ll (Nu)^\eta |L(s,\psi)|V^{c+d}\log^2(2V)$, with constants
retaining fixed margins, $\nu,S$, and $\sigma$. Choosing $c,d$ close to
$1-\sigma$ reproduces $U^{2-2\sigma+\eta}$ in amplitude and
$U^{2-4\delta+\eta}$ in energy on $\sigma=1/2+\delta$.
Insertion into the original common-profile Mellin argument retains
$M_4(W)^2$ and all previously required height/conductor uniformities.
No derivative-free improvement follows from this Perron estimate.

The correction converges even on some negative auxiliary lines, but
moving there can cross $z=\rho-s$ or $w=\rho-s$, from zeros of the
primitive $L$-function. The exact meromorphic slice
$\mathcal F(s,z,0)=1/L^{T_0}(s+z,\psi)$ forbids assuming universal
cancellation of these poles. At a numerator zero the slice must be
understood meromorphically, not by an unjustified pointwise substitution.
Sequential shifts have single boundary integrals and crossed-zero
residues as well as the kernel double residue; away from $s$ zeros or
poles that double residue is $1/L^{T_0}(s,\psi)$.
A principal $L$ pole supplies a reciprocal zero, not a reciprocal pole.
Uniform control of actual reciprocal integrals/residues or the signed
covariance remains a new input. Numerator AFM alone supplies neither.

% Generated shared reviewed material.
\section{A quantitative covariance budget without a new exponent}\label{cbq:core}
Retain the original Eisenstein sextic symbols, fixed finite-order $\nu$,
fixed $S$, one possibly complex common $W\in C_c^\infty$ supported in
$[1,2]$, all squarefree columns prime to $S$, and every raw shared-prime
zero. In the accepted window $2D/7\le U^2\le D$, put $r_m=r_U(m)$,
$\lambda_m=\mu(m)\nu(m)W(Nm/D)$ and $v_m=\overline{\lambda_m}$.
Thus $C_U=\mu r$ and $A_U=\mu(r-1)$ on the active annulus.
For a finite literal row block $B$, let
\[
 G_{mn}=\sum_{u\in B}x_m(u)\overline{x_n(u)},\qquad
 R=\operatorname{diag}(r_m),\qquad
 M_j=\sum_mr_m^j|W(Nm/D)|^2\quad(0\le j\le2).
\]
The primitive $\psi_u,\psi_{u,2}$ and full deleted sets retain their
previous meanings; no conductor-drop value replaces a raw zero.

\begin{proposition}[An actual-vector lower budget]\label{cbq:budget}
If $0\le G\le\Lambda_B I$, then
\begin{equation}\label{cbq:vector}
 \Re Q_B(C_U,\mu)\ge
 -\frac{\Lambda_B}{2D}\{\sqrt{M_0M_2}-M_1\}.
\end{equation}
In this high window, uniformly in $U$,
\begin{equation}\label{cbq:moments}
 \sum_{D\le Nm\le2D}r_m\ll_K D,\qquad
 \sum_{D\le Nm\le2D}r_m^2\ll_K D\log(2D).
\end{equation}
Consequently
\[
 \Re Q_B(C_U,\mu)\ge
 -C_K\|W\|_\infty^2\Lambda_B\sqrt{\log(2D)}.
\]
\end{proposition}
The vector inequality is sharp among positive contractions with the
specified $v,Rv$; no such sharpness is asserted for the arithmetic Gram.
The moment bound counts actual balanced ideal factorizations, rather than
assuming independence of the two row factors.

On $B_X$, inserting the previously accepted powered-BGL norm gives only
\[
 \Re Q_{B_X}(C_U,\mu)\ge
 -C_{\nu,S,\epsilon}\|W\|_\infty^2(DX)^\epsilon
 \{X^2+DX+D^{2/3}X^{5/3}\}.
\]
At $X\asymp D^{1/4}$ the exponents remain $1/2,5/4,13/12$.
At $H=D^{3/4}$ this budget still has the old $D^{1/2}$ excess over $H$;
it supplies no $BT_\sigma$ saving. A negative-part budget is weaker than
$2\Re Q(C_U,\mu)\ge Q(C_U,C_U)+Q(\mu,\mu)-R_\sigma$.

For $S_R=\sqrt R$, the exact identity
\[
 \tfrac12(RG+GR)=S_RGS_R+\tfrac12[S_R,[S_R,G]]
\]
localizes the commutator to unequal factor-count levels, with entry
$(\sqrt{r_m}-\sqrt{r_n})^2G_{mn}$. If $B$ retains all six literal units,
the actual unit sum further gives $G_{mn}=0$ unless $Nm\equiv Nn\pmod{36}$.
These are exact deletions, not a power saving. Even zero different-level
pairings leave the nonnegative within-level energies
$\sum_k(k-1)^2Q_B(\mu\mathbf1_{r=k},\mu\mathbf1_{r=k})$.
The unequal-product coupled fourth moment and its mixed M\"obius term
remain unestimated. No new analytic exponent, energy lower bound,
Lean formalization or GRH conclusion follows.

% Generated shared reviewed material.
\section{Proof and arithmetic boundary of the covariance budget}\label{cbq:details}
\subsection{The positive-contraction calculation}
Since $G$ is the exact raw Gram, it is positive semidefinite. If
$\Lambda_B>0$, put $T=2G/\Lambda_B-I$, a self-adjoint contraction. Then
\[
 \Re(v^*RGv)=\frac{\Lambda_B}{2}
 \{M_1+\Re(v^*RTv)\}
 \ge\frac{\Lambda_B}{2}\{M_1-\|Rv\|_2\|v\|_2\},
\]
which is~\eqref{cbq:vector}; $\Lambda_B=0$ is immediate.
The Hermitian rank-two matrix
$\{(Rv)v^*+v(Rv)^*\}/2$ has possible nonzero eigenvalues
$(M_1\pm\sqrt{M_0M_2})/2$. Its negative spectral projection realizes
the lower bound among all $0\le G\le\Lambda_B I$, a linear-algebra
sharpness statement only. If $a\le r_m\le b$ on the weighted support,
$a+b>0$, the pointwise inequality $r_m^2\le(a+b)r_m-ab$ also gives
\[
 \Re Q_B(C_U,\mu)\ge
 -\frac{\Lambda_B M_0(b-a)^2}{8D(a+b)}.
\]
For $a=b=0$ the covariance is zero. The constant follows by maximizing
$\sqrt{(a+b)t-ab}-t$ on $[a,b]$.

\subsection{The balanced-divisor second moment}
Set $T_0=\sqrt D$. The high-window factor restrictions imply
\[
 \sqrt{2/7}\,T_0<Na,Nb<\sqrt{14}\,T_0.
\]
The elementary ideal count $\#\{a:Na\le Y\}\ll_K Y$ proves the
first assertion of~\eqref{cbq:moments}. For the second, two ordered
factorizations $ab=cd=m$ have the unique form
\[
 a=gh,\quad c=gk,\quad b=k\ell,\quad d=h\ell,
 \qquad g=(a,c).
\]
Squarefreeness makes $g,h,k,\ell$ pairwise coprime and squarefree.
To upper-bound this positive count, enlarge to arbitrary ideals.
For $M=\max(Nh,Nk)$ the individual bounds are
\[
 Ng\le\frac{\sqrt{14}\,T_0}{M},\qquad
 N\ell\le\frac{\sqrt{14}\,T_0}{M},\qquad
 1\le M\le\sqrt{14}\,T_0.
\]
On $2^j\le M<2^{j+1}$ there are $O_K(2^{2j})$ choices of $(h,k)$
and $O_K((T_0/2^j)^2)$ choices of $(g,\ell)$, with bounded terminal
ranges absorbed in the constant. Each dyadic range contributes $O_K(D)$;
there are $O(\log(2D))$ ranges. This proves the second moment with no
extra logarithm from counting only $NgN\ell$. Restricting to primes
outside fixed $S$ decreases the count. Hence
\[
 M_0,M_1\ll_K\|W\|_\infty^2D,\qquad
 M_2\ll_K\|W\|_\infty^2D\log(2D).
\]
These constants are independent of the cutoff in the stated window.
Only the common supremum norm is used, including for complex $W$.

\subsection{What the commutator actually controls}
Write
\[
 t_{mn}=D^{-1}\Re\{\lambda_m\overline{\lambda_n}G_{mn}\},\qquad
 \delta_{mn}=(\sqrt{r_m}-\sqrt{r_n})^2.
\]
The exact commutator identity proves
\[
 \Re Q_B(C_U,\mu)\ge
 Q_B(\sqrt r\,\mu,\sqrt r\,\mu)
 -\tfrac12\sum_{m,n}\delta_{mn}|t_{mn}|.
\]
The displayed loss retains the actual character and profile phases before
taking absolute values; it is at most the phase-blind loss obtained from
$|t_{mn}|\le D^{-1}|\lambda_m\lambda_nG_{mn}|$.
On the accepted four-column smooth $B_7$ fixture at amplitude $x\ge0$,
the positive term is $18/25$, this coefficient-adapted loss is $9x/50$,
and the lower bound is attained. Purely imaginary entries of the full
complex Gram make the phase-blind loss larger. This calibrates the new
budget and does not give an asymptotic obstruction.

For the unit deletion let $\zeta_6=1+\omega$. Every good prime norm $q$
is $1\pmod6$, and direct residue-field exponentiation gives
$x_p(\zeta_6)=\zeta_6^{(q-1)/6}$. Multiplying over $p\mid m$ and reducing
$Nm=\prod_{p\mid m}(1+6e_p)$ modulo $36$ gives
\[
 x_m(\zeta_6)=\zeta_6^{(Nm-1)/6\bmod6}.
\]
Multiplication of all literal rows by $\zeta_6$ preserves their raw zeros
and proves the norm-congruence deletion in the core. The vector budget
may be applied separately on these six classes. No D-dependent number
of classes or conductor saving is produced.

The moment proof above counts the coefficient diagonal $ab=cd$.
It has not estimated the unequal-product character quadruples in the
coupled row fourth moment. Replacing the inherited $DX$ norm by
$D^{1-\sigma}X$ merely because equal-level and different-class entries
vanish is the unsupported step. Small negative commutator loss alone
also leaves the within-level energies in the core. Principal residues
are unchanged; this covariance is used on the nonprincipal blocks where
the original residue is zero, not as a new centering prescription.

\subsection{The complete near-saturating face}
For a six-factor row block let $R_0=\prod_{j=1}^6A_j$,
$M=\max(A_1,A_2)$ and $\prod_j A_j^j\le H$.
The accepted exact-mask block norm is
\[
 \Lambda_A\ll_{\nu,S,\epsilon}(HD)^\epsilon
 \{R_0+DR_0/M+D^{2/3}R_0/M^{1/3}\}.
\]
The column moment and covariance proofs apply to every such block.
Write $x_j=\log A_j$, $h=\log H$, $c_0=\sum_jjx_j$ and
$r_0=\sum_jx_j$. Then, for $x_1\ge x_2$,
\[
 h-3(r_0-\max(x_1,x_2))
 =h-c_0+(x_1-x_2)+x_4+2x_5+3x_6;
\]
for $x_2\ge x_1$ replace $x_1-x_2$ by $2(x_2-x_1)$.
All terms are nonnegative. Thus the middle monomial is saturated on
\[
 A_1=A_2,\qquad A_1A_3=H^{1/3},\qquad A_4=A_5=A_6=1,
\]
with full capacity. The balanced $A_3=1$ case is only an endpoint.
The $a_3$ direction remains nonprincipal by its outside-$S$ valuation
$3\pmod6$. If $R_0/M\ge H^{1/3}D^{-\eta}$, the total displayed slack
is at most $3\eta\log D$. Therefore a fixed power gain must be uniform
on this full neighboring region. The proved budget still costs $DR_0/M$;
it gives no global improvement over the accepted envelope.

The independent standard-library diagnostic checks nine ideal moment
cases, eight raw Gram cases with $27\,000$ phase values and $5\,184$ raw
zeros, exact radical commutator/vector budgets, and $227$ exponent cases.
Finite tables use common norm profiles realizable by smooth disjoint bumps;
they are not tests of an infinite analytic estimate. Full derivation,
source bindings and preserved checks are under
\path{planning/covariance-budget-20261105/attempt/} and
\path{outputs/checks/2026-11-05-covariance-budget/attempt/}.

% Generated shared reviewed material.
\section{The fourth-moment conductor and its full raw zero set}
\label{upq:core}
Retain the original Eisenstein sextic symbols $x_m(u)=(u/m)_6$,
fixed finite-order $\nu$, fixed $S$, and one possibly complex common
smooth $W$ supported in $[1,2]$. Put
$\lambda_m=\mu(m)\nu(m)W(Nm/D)$ and $r_m=r_U(m)$.
Columns remain squarefree and prime to $S$. In $2D/7\le U^2\le D$, the actual polynomial is
\[
 \mathcal C_U(u)=
 \sum_{\substack{(a,b)=1,\ Na,Nb>U\\D\le Na\,Nb\le2D}}
 \mu(a)\mu(b)\nu(ab)W(Nab/D)x_a(u)x_b(u).
\]
Its squared row energy has four factor indices $a,b,c,d$, with
$(a,b)=(c,d)=1$. These internal masks are essential. At each prime,
\[
 e_p=\mathbf1_{p\mid a}+\mathbf1_{p\mid b}
       -\mathbf1_{p\mid c}-\mathbf1_{p\mid d}
 \quad\hbox{belongs to }\{0,\pm1\}.
\]
Thus, for the two squarefree products $m=ab,n=cd$, the primitive
oscillatory conductor outside fixed $S$ is $q=mn/(m,n)^2$.
The full raw zero-preserving period is $L=mn/(m,n)$:
\[
 x_m(u)\overline{x_n(u)}
 =\mathbf1_{(u,(m,n))=1}\theta_{m/(m,n),\,n/(m,n)}(u).
\]
No common-prime zero is restored by taking the primitive character.

A generic fourth moment $|\sum_m\lambda_mx_m(u)|^4$ has four individually
squarefree columns $m_1,\ldots,m_4$ with no such internal coprimality.
Its signed local exponent
$e_p=\sum_{i=1}^2\mathbf1_{p\mid m_i}-\sum_{i=3}^4\mathbf1_{p\mid m_i}$
can also be $\pm2$, giving a nontrivial cubic component. For $L=\operatorname{rad}(m_1m_2m_3m_4)$, the exact row-power
interface is
\[
 Q_j=\prod_{\substack{p\mid L\\je_p\not\equiv0\ (6)}}p,\qquad
 Z_j=L/Q_j,\qquad
 \prod_{i=1}^2x_{m_i}(v^j)
       \overline{\prod_{i=3}^4x_{m_i}(v^j)}
 =\theta_{j,e}(v)\mathbf1_{(v,Z_j)=1}.
\]
Here $\theta_{j,e}$ is primitive on $Q_j$, with its own zero extension.
At $j=3$ an exponent $\pm2$ becomes primitive-principal but its raw
zero remains. For the original two-squarefree-product expansion,
$Q_j=q$ for all $1\le j\le5$, including the quadratic $j=3$ direction.
For the accepted row blocks $Nr_j\asymp A_j$, put
$R_0=\prod_jA_j$, $P=\max(A_1,\ldots,A_5)$ and retain
$\prod_jA_j^j\le H$. The full $A_3$ saturation face and all row-factor
masks remain.

The small-conductor pair count recovers an already controlled high-gcd
sector; it is not a new analytic advance. On the remaining original
sector $Nq>P^2$, favorable smooth completion
costs $\min(A_j,\sqrt{Nq})=A_j$ for $1\le j\le5$.
The literal sharp cutoff is not replaced by this smooth trial.
The signed large-conductor fourth-moment and mixed M\"obius sums remain
open; no improved full exponent, $BT_\sigma$, new Lean theorem,
zero-location result or GRH conclusion follows.

% Generated shared reviewed material.
\section{Scope of the raw fourth-moment reduction}
\label{upq:details}
Write $c_{a,b}=\mu(a)\mu(b)\nu(ab)W(Nab/D)$ on the exact factor support
in the core. Before absolute Cauchy, the normalized energy is
\[
 \frac1D\sum_{\substack{(a,b),(c,d)\ {\rm on\ that\ support}}}
 c_{a,b}\overline{c_{c,d}}
 \sum_{u\in B}x_a(u)x_b(u)\overline{x_c(u)x_d(u)}.
\]
Squarefreeness of $ab$ and $cd$ forces each positive and negative
prime multiplicity to be at most one, proving $e_p\in\{0,\pm1\}$.
For the mixed term against the actual $\mu$ polynomial the exponent is
$\mathbf1_{p\mid a}+\mathbf1_{p\mid b}-\mathbf1_{p\mid n}$,
with the same conclusion and the same full union mask.
The common possibly complex $W$, $\nu$, hyperbolic support and all raw
zeros stay in these identities.

In contrast, unrestricted four-column memberships $(1,1,0,0)$ give
$e_p=2$, a cubic component; $(1,1,1,0)$ give $e_p=1$ although both
product supports contain $p$. Neither the parity symmetric difference
of four supports nor the symmetric difference of two nonsquarefree
product supports gives the general conductor.
For the interface in the core, its local character order is
\[
 \begin{array}{c|cccc}
 j&1,5&2,4&3&6\\ \hline
 e_p=\pm1&6&3&2&1\\
 e_p=\pm2&3&3&1&1
 \end{array}
\]
Nontrivial entries have tame conductor exponent one; trivial entries
remain deleted at $p\mid v$. Multiplying the interface over all six
row factors, with the fixed unit/$S$ phase, preserves every original
zero, including those at sixth-power factors. All inter-factor,
squarefree, ray and capacity restrictions remain external restrictions
on the same row sum; this identity asserts no independence.

At a residue field of size $Q$, let $\chi$ have order six and $\rho$
be nontrivial additive. The zero-extended complete sum
$G_k(t)=\sum_{z\ne0}\chi(z)^k\rho(tz)$ satisfies
\[
 \begin{array}{c|cc}
 &t=0&t\ne0\\ \hline
 k\not\equiv0\ (6)&0&|G_k(t)|=\sqrt Q\\
 k\equiv0\ (6)&Q-1&-1.
 \end{array}
\]
For $t\ne0$, substitution gives
$G_k(t)=\overline{\chi(t)^k}G_k(1)$; for nontrivial $k$,
expanding the norm square gives
$(Q-1)-\sum_{a\ne0,1}\chi(a)^k=Q$.
The proof applies at split and inert primes. By CRT the full period-$L$
transform vanishes if any local frequency on $Q_j$ is zero; otherwise
its magnitude is
\[
 (NQ_j)^{1/2}
 \prod_{p\mid Z_j}
 \left|-1+Np\,\mathbf1_{\text{local frequency at }p\text{ is }0}\right|.
\]
Canceled common primes thus supply Ramanujan factors.

Exact deletion, with all remaining restrictions kept in $F$, gives
\[
 \sum_r\theta(r)F(r)\mathbf1_{(r,Z)=1}
 =\sum_{d\mid Z}\mu(d)\theta(d)\sum_z\theta(z)F(dz).
\]
For a squarefree pivot let $E$ be the radical of its fixed excluded
primes, including $Z$ and the relevant other row factors, with
conductor primes omitted since $\theta$ already vanishes there.
The additional squarefree identity is
\[
 \mu(r)^2\mathbf1_{(r,E)=1}
 =\sum_{\substack{h^2\mid r\\(h,E)=1}}\mu(h)
       \sum_{\substack{d\mid E\\d\mid r}}\mu(d).
\]
Since $(d,h)=1$, its lifted row sum has terms
$\mu(d)\mu(h)\theta(d)\theta(h)^2\sum_z\theta(z)F(dh^2z)$.
The scaled $F$ retains the exact window, primary/ray restrictions and
other masks. Divisor triangles cost at most $(HD)^\epsilon$;
they do not restore a forbidden nonzero.

The previous fixed-radial-profile high-gcd proof
(\path{research/short-row-high-gcd.md}, independently reviewed in
\path{planning/reviews/book/short-row-formal.md})
already bounds its off-diagonal tail by
$O_{\Phi,\epsilon}(D^{2+\epsilon}/G_0^2)$ for $Nh\ge G_0$.
For annular $m=hr,n=ht$, $Nq\le Y$ implies $Nh\ge D/\sqrt Y$,
so the old smooth-scope bound is $O_{\Phi,\epsilon}(D^\epsilon Y)$.
The present elementary count $O(D\sqrt Y)$, factor multiplicities
$r_U(m)r_U(n)\ll_\epsilon D^\epsilon$, and $\#B_A=O(R_0)$ give only
$O(\|W\|_\infty^2D^\epsilon R_0\sqrt Y)$ on a literal sharp block.
This is a weaker conductor-form recovery in its own bookkeeping scope,
not a new high-gcd theorem; the smooth estimate is not silently
transferred to a selected sharp off-diagonal piece.
The full $A_3$ face and neighborhood slack are also inherited.

Finally, primitive smooth Poisson completion has bound
$O_V(\min(A_j,\sqrt{Nq}))$ on the original support.
For $Nq>P^2$ and $1\le j\le5$, every scaled length
$A_j/(Nd\,Nh^2)$ is below $\sqrt{Nq}$, so even this favorable trial
takes its trivial length term. Its constant uses fixed smooth
seminorms and need not stay bounded when a sharp boundary is approached.
In the admissible subsector $(m,n)=1$, $Nq=Nm\,Nn\asymp D^2$.
A complete zero mean does not control the nonzero frequencies of the
short original masked row window. The unresolved target is their
joint signed sum with weights $r_U(m)r_U(n)\lambda_m\overline{\lambda_n}$
and $r_U(m)\lambda_m\overline{\lambda_n}$, before per-conductor absolute
estimates. No new global analytic conclusion is claimed.


\end{document}
